% Complete reviewed Unit 13; editorial master: chapters/ch13/answer-key.md
% Markdown SHA-256: 8f997622ba66ec27326bca1fd9c88174064046eda3b236f5ec10e2e035a19811
% One-use initial assembly followed by complete content and page review.
\newcommand{\runningtitle}{Unit 13. Answer key}
% Layout inherited from the reviewed previous course; set a running title per unit.
% This is a small hand-maintained TeX layout, not a conversion pipeline.
\RequirePackage{fix-cm}
\documentclass[11pt,twoside]{article}
\usepackage{fontspec}
\setmainfont{Cambria}[
  Renderer=HarfBuzz,
  BoldFont=cambriab.ttf,
  ItalicFont=cambriai.ttf,
  BoldItalicFont=cambriaz.ttf
]
\setmonofont{lmmono10-regular.otf}[BoldFont=lmmonolt10-bold.otf]
\usepackage[paperwidth=165mm,paperheight=241mm,inner=21mm,outer=19mm,
  top=22mm,bottom=20mm,headheight=14pt,headsep=7mm,footskip=10mm]{geometry}
\usepackage{xcolor,array,booktabs,longtable,tabularx}
\usepackage{enumitem,fancyhdr}
\usepackage[unicode,hidelinks]{hyperref}
\definecolor{ink}{HTML}{000000}
\definecolor{muted}{HTML}{4C5558}
\setlength{\parindent}{0pt}
\setlength{\parskip}{3.5pt}
\AtBeginDocument{\fontsize{11}{14.3}\selectfont}
\setlength{\tabcolsep}{4pt}
\renewcommand{\arraystretch}{1.10}
\setlist[itemize]{leftmargin=1.3em,label=\textbullet,itemsep=2pt,topsep=3pt,parsep=0pt}
\setlist[enumerate]{leftmargin=1.8em,itemsep=3pt,topsep=4pt,parsep=0pt}
\widowpenalty=10000
\clubpenalty=10000
\emergencystretch=1em
\raggedbottom
\pagestyle{fancy}
\fancyhf{}
\providecommand{\runningtitle}{Old English}
\fancyhead[LE,RO]{\small\itshape \runningtitle}
\fancyhead[LO,RE]{\small Old English}
\fancyfoot[LE,RO]{\small\thepage}
\renewcommand{\headrulewidth}{0.3pt}
\newcommand{\booktitle}[2]{
  {\fontsize{9}{11}\selectfont\color{muted}#1\par}
  \vspace{3pt}
  {\fontsize{20}{24}\selectfont\color{ink}#2\par}
  \vspace{6pt}
}
\newcommand{\source}[1]{\par{\footnotesize\color{muted}#1\par}}
\newcommand{\readingtext}[1]{
  \par\addvspace{5pt}
  \begingroup\leftskip=1.5em
    \fontsize{12}{16}\selectfont #1\par
  \endgroup\addvspace{5pt}
}
\makeatletter
\renewcommand{\section}{\@startsection{section}{1}{\z@}
  {-2.4ex \@plus -0.6ex \@minus -0.2ex}{0.8ex \@plus 0.2ex}
  {\normalfont\fontsize{12}{15}\selectfont\bfseries\color{ink}}}
\renewcommand{\subsection}{\@startsection{subsection}{2}{\z@}
  {-1.8ex \@plus -0.5ex \@minus -0.2ex}{0.5ex \@plus 0.2ex}
  {\normalfont\fontsize{11}{14}\selectfont\bfseries\color{ink}}}
\makeatother

\hyphenpenalty=10000
\exhyphenpenalty=10000
\emergencystretch=3em
\newcommand{\itempara}[2]{\par\addvspace{3pt}\noindent\hangindent=1.8em\hangafter=1\makebox[1.8em][l]{#1.}#2\par}
\newcommand{\reserve}[1]{\par\penalty-100\begingroup\dimen0=\pagegoal\advance\dimen0 by-\pagetotal\ifdim\dimen0<#1\newpage\fi\endgroup}
\begin{document}
\addfontfeatures{Ligatures=TeXOff}
\raggedright
\booktitle{UNIT 13}{Answer key}

% mdparagraph:P001:begin
This key covers all 90 tasks: the four complete Cook sets and 38 original exercises. The models, translations and explanations are new teaching material. The twelve translated examples are teaching examples, not additional exercises. Equivalent English wording and defensible Old English alternatives are acceptable. A grammatical answer can nevertheless miss an explicit target, such as object-first order, a particular mood or the selected paradigm. [1]
% mdparagraph:P001:end

% mdparagraph:P002:begin
In a form task, give every labelled cell, including cells whose spellings coincide. In sentence work, preserve the participants, case relations, agreement, clause boundaries and intended reference. Do not identify subjects by first position alone, count a multiword constituent as several constituents, or choose a pronoun’s antecedent merely because it is the nearest noun. A masculine Old English pronoun can refer to a masculine noun meaning ‘song’.
% mdparagraph:P002:end

% mdparagraph:P003:begin
The local Cook and reading helps remain available. Their supplied forms and phrases are not tests of unaided recall. A complete indicative table does not assign its verb’s infinitive, commands or participial constructions; use the separately taught or supplied systems. In translation, ordinary preterites may express a prior event when the context establishes it. Do not insert an auxiliary merely because the English prompt has ‘shall’, ‘was going’ or ‘had come’.
% mdparagraph:P003:end

\clearpage
% mdparagraph:P004:begin
\section*{F. Complete forms}
% mdparagraph:P004:end

\reserve{5\baselineskip}
\reserve{8\baselineskip}
% mdparagraph:P005:begin
\textbf{F01. The three kinship nouns.}
% mdparagraph:P005:end

% item: NEW-U13-F01
% key: KEY-NEW-U13-F01
\reserve{9\baselineskip}
% mdparagraph:P006:begin
\textbf{Mōdor}
% mdparagraph:P006:end

% mdtable:T01:begin
\par\begingroup
\fontsize{10.5}{13.2}\selectfont
\setlength{\LTleft}{0pt}\setlength{\LTright}{\fill}
\noindent
\begin{tabular}{@{}>{\raggedright\arraybackslash}p{39.46mm}>{\raggedright\arraybackslash}p{39.46mm}>{\raggedright\arraybackslash}p{39.46mm}@{}}
\toprule
\textbf{Case} & \textbf{Singular} & \textbf{Plural} \\
\midrule
Nominative & mōdor & mōdru \\
Genitive & mōdor & mōdra \\
Dative & mēder & mōdrum \\
Accusative & mōdor & mōdru \\
\bottomrule
\end{tabular}
\par\endgroup

% mdtable:T01:end
\reserve{9\baselineskip}
% mdparagraph:P007:begin
\textbf{Dohtor}
% mdparagraph:P007:end

% mdtable:T02:begin
\par\begingroup
\fontsize{10.5}{13.2}\selectfont
\setlength{\LTleft}{0pt}\setlength{\LTright}{\fill}
\noindent
\begin{tabular}{@{}>{\raggedright\arraybackslash}p{39.46mm}>{\raggedright\arraybackslash}p{39.46mm}>{\raggedright\arraybackslash}p{39.46mm}@{}}
\toprule
\textbf{Case} & \textbf{Singular} & \textbf{Plural} \\
\midrule
Nominative & dohtor & dohtor \\
Genitive & dohtor & dohtra \\
Dative & dehter & dohtrum \\
Accusative & dohtor & dohtor \\
\bottomrule
\end{tabular}
\par\endgroup

% mdtable:T02:end
\reserve{9\baselineskip}
% mdparagraph:P008:begin
\textbf{Brōðor}
% mdparagraph:P008:end

% mdtable:T03:begin
\par\begingroup
\fontsize{10.5}{13.2}\selectfont
\setlength{\LTleft}{0pt}\setlength{\LTright}{\fill}
\noindent
\begin{tabular}{@{}>{\raggedright\arraybackslash}p{39.46mm}>{\raggedright\arraybackslash}p{39.46mm}>{\raggedright\arraybackslash}p{39.46mm}@{}}
\toprule
\textbf{Case} & \textbf{Singular} & \textbf{Plural} \\
\midrule
Nominative & brōðor & brōðor \\
Genitive & brōðor & brōðra \\
Dative & brēðer & brōðrum \\
Accusative & brōðor & brōðor \\
\bottomrule
\end{tabular}
\par\endgroup

% mdtable:T03:end
% mdparagraph:P009:begin
\textbf{Explanation:} The selected genitive singular equals the nominative in all three nouns. The dative singular changes the vowel: \textbf{mēder, dehter, brēðer}. Before plural \textbf{-a, -um}, use the syncopated stems \textbf{mōdr-, dohtr-, brōðr-}. The selected nominative/\allowbreak{}accusative plurals are \textbf{mōdru, dohtor, brōðor}. Preserve short \textbf{o} in \textbf{dohtor}; the long vowels of the other nouns do not establish its quantity. Documented nominative/\allowbreak{}accusative plural alternatives \textbf{mōdra, dohtru/\allowbreak{}dohtra, brōðru} are valid in those cells. The later genitives \textbf{mēder, dehter} are documented variants, but the displayed table gives the requested selection. Extending a mutated dative stem through every plural cell misses the pattern. [2]
% mdparagraph:P009:end

\reserve{5\baselineskip}
\par\begin{minipage}{\linewidth}\raggedright
% mdparagraph:P010:begin
\textbf{F02. Æx and ēagðȳrel.}
% mdparagraph:P010:end

% item: NEW-U13-F02
% key: KEY-NEW-U13-F02
% mdparagraph:P011:begin
\textbf{Æx}
% mdparagraph:P011:end

% mdtable:T04:begin
\par\begingroup
\fontsize{10.5}{13.2}\selectfont
\setlength{\LTleft}{0pt}\setlength{\LTright}{\fill}
\noindent
\begin{tabular}{@{}>{\raggedright\arraybackslash}p{39.46mm}>{\raggedright\arraybackslash}p{39.46mm}>{\raggedright\arraybackslash}p{39.46mm}@{}}
\toprule
\textbf{Case} & \textbf{Singular} & \textbf{Plural} \\
\midrule
Nominative & æx & æxa \\
Genitive & æxe & æxa \\
Dative & æxe & æxum \\
Accusative & æxe & æxa \\
\bottomrule
\end{tabular}
\par\endgroup

% mdtable:T04:end
% mdparagraph:P012:begin
\textbf{Ēagðȳrel}
% mdparagraph:P012:end

% mdtable:T05:begin
\par\begingroup
\fontsize{10.5}{13.2}\selectfont
\setlength{\LTleft}{0pt}\setlength{\LTright}{\fill}
\noindent
\begin{tabular}{@{}>{\raggedright\arraybackslash}p{39.46mm}>{\raggedright\arraybackslash}p{39.46mm}>{\raggedright\arraybackslash}p{39.46mm}@{}}
\toprule
\textbf{Case} & \textbf{Singular} & \textbf{Plural} \\
\midrule
Nominative & ēagðȳrel & ēagðȳrlu \\
Genitive & ēagðȳrles & ēagðȳrla \\
Dative & ēagðȳrle & ēagðȳrlum \\
Accusative & ēagðȳrel & ēagðȳrlu \\
\bottomrule
\end{tabular}
\par\endgroup

% mdtable:T05:end
% mdparagraph:P013:begin
\textbf{Explanation:} Feminine \textbf{æx} keeps \textbf{æx-} before its endings; the singular oblique cases take \textbf{-e}. Do not import the masculine short-stem \textbf{æ/\allowbreak{}a} alternation. Neuter \textbf{ēagðȳrel} keeps the full endingless nominative/\allowbreak{}accusative singular, but drops the unstressed vowel before an inflectional ending: \textbf{ēagðȳrl-es, -e, -u, -a, -um}. The plural \textbf{ēagðȳrlu} is not obtained by adding \textbf{-u} to the entire singular. Preserve the long \textbf{ȳ}. [2]
% mdparagraph:P013:end

\end{minipage}\par\addvspace{4pt}

\par\begin{minipage}{\linewidth}\raggedright
% mdparagraph:P014:begin
\textbf{F03. Three diagnostic patterns.}
% mdparagraph:P014:end

% item: NEW-U13-F03
% key: KEY-NEW-U13-F03
% mdtable:T06:begin
\par\begingroup
\fontsize{10.5}{13.2}\selectfont
\setlength{\LTleft}{0pt}\setlength{\LTright}{\fill}
\noindent
\begin{tabular}{@{}>{\raggedright\arraybackslash}p{18.38mm}>{\raggedright\arraybackslash}p{24.17mm}>{\raggedright\arraybackslash}p{22.51mm}>{\raggedright\arraybackslash}p{24.17mm}>{\raggedright\arraybackslash}p{23.52mm}@{}}
\toprule
\textbf{Noun} & \textbf{Genitive singular} & \textbf{Dative singular} & \textbf{Nominative plural} & \textbf{Genitive plural} \\
\midrule
cirice & cirican & cirican & cirican & ciricena \\
mynster & mynstres & mynstre & mynstru & mynstra \\
wyrttruma & wyrttruman & wyrttruman & wyrttruman & wyrttrumena \\
\bottomrule
\end{tabular}
\par\endgroup

% mdtable:T06:end
% mdparagraph:P015:begin
\textbf{Explanation:} \textbf{Cirice} follows the weak feminine model, using \textbf{ciric-} before \textbf{-an} and \textbf{-ena}. Its selected ordinary forms remain distinct from the reading’s supplied \textbf{ciricean} spellings. \textbf{Mynster} is the selected syncopating neuter: \textbf{mynstr-} before \textbf{-es, -e, -u, -a}. \textbf{Wyrttruma} follows the weak masculine model: \textbf{wyrttrum-} before \textbf{-an} and genitive plural \textbf{-ena}. A masculine \textbf{-as} plural or an unsyncopated selected \textbf{mynsteru} would miss the requested models. [2]
% mdparagraph:P015:end

\end{minipage}\par\addvspace{4pt}

\reserve{8\baselineskip}
% mdparagraph:P016:begin
\textbf{F04. Swēte and its comparison forms.}
% mdparagraph:P016:end

% item: NEW-U13-F04
% key: KEY-NEW-U13-F04
% mdtable:T07:begin
% course-table-style: 2026-10
% course-panel:T07-P01:begin
\par\begingroup
\fontsize{10.5}{13.2}\selectfont
\setlength{\LTleft}{0pt}\setlength{\LTright}{\fill}
\noindent
\begin{tabular}{@{}>{\raggedright\arraybackslash}p{22.0mm}>{\raggedright\arraybackslash}p{31.188mm}>{\raggedright\arraybackslash}p{31.188mm}>{\raggedright\arraybackslash}p{31.188mm}@{}}
\toprule
\textbf{Case} & \textbf{Masculine singular} & \textbf{Neuter singular} & \textbf{Feminine singular} \\
\midrule
Nominative & swēte & swēte & swētu \\
Genitive & swētes & swētes & swētre \\
Dative & swētum & swētum & swētre \\
Accusative & swētne & swēte & swēte \\
\bottomrule
\end{tabular}
\par\endgroup
% course-panel:T07-P01:end
\par\addvspace{5pt}
% course-panel:T07-P02:begin
\par\begingroup
\fontsize{10.5}{13.2}\selectfont
\setlength{\LTleft}{0pt}\setlength{\LTright}{\fill}
\noindent
\begin{tabular}{@{}>{\raggedright\arraybackslash}p{22.0mm}>{\raggedright\arraybackslash}p{31.188mm}>{\raggedright\arraybackslash}p{31.188mm}>{\raggedright\arraybackslash}p{31.188mm}@{}}
\toprule
\textbf{Case} & \textbf{Masculine plural} & \textbf{Neuter plural} & \textbf{Feminine plural} \\
\midrule
Nominative & swēte & swētu & swēta \\
Genitive & swētra & swētra & swētra \\
Dative & swētum & swētum & swētum \\
Accusative & swēte & swētu & swēta \\
\bottomrule
\end{tabular}
\par\endgroup
% course-panel:T07-P02:end
% mdtable:T07:end

% mdtable:T08:begin
\par\begingroup
\fontsize{10.5}{13.2}\selectfont
\setlength{\LTleft}{0pt}\setlength{\LTright}{\fill}
\noindent
\begin{tabular}{@{}>{\raggedright\arraybackslash}p{89.83mm}>{\raggedright\arraybackslash}p{31.36mm}@{}}
\toprule
\textbf{Requested form} & \textbf{Answer} \\
\midrule
Comparative, weak masculine nominative singular & swētra \\
Superlative, strong masculine nominative singular & swētest \\
Superlative, weak masculine nominative singular & swētesta \\
\bottomrule
\end{tabular}
\par\endgroup

% mdtable:T08:end
% mdparagraph:P017:begin
\textbf{Explanation:} Positive \textbf{swēte} preserves \textbf{-e} in its selected endingless strong cells; the other forms use \textbf{swēt-} and the displayed strong endings. Compare neuter singular \textbf{swēte} with neuter plural \textbf{swētu}. The comparative uses weak inflection only: its masculine nominative singular is \textbf{swētra}. The superlative stem is \textbf{swētest-}: strong masculine nominative singular \textbf{swētest}, weak \textbf{swētesta}. Comparison changes the stem; strong or weak selection supplies the appropriate ending. Do not invent a strong comparative paradigm or use weak accusative \textbf{swētestan} as a nominative. The full comparative’s genitive plural is \textbf{swētra}, with the checked consonant simplification, not a mechanically doubled \textbf{rr}. [2]
% mdparagraph:P017:end

\reserve{11\baselineskip}
\par\begin{minipage}{\linewidth}\raggedright
% mdparagraph:P018:begin
\textbf{F05. Forceorfan: complete indicative.}
% mdparagraph:P018:end

% item: NEW-U13-F05
% key: KEY-NEW-U13-F05
% mdtable:T09:begin
% course-table-style: 2026-10
% course-panel:T09-P01:begin
\par\begingroup
\fontsize{10.5}{13.2}\selectfont
\setlength{\LTleft}{0pt}\setlength{\LTright}{\fill}
\noindent
\begin{tabular}{@{}>{\raggedright\arraybackslash}p{39.459mm}>{\raggedright\arraybackslash}p{39.459mm}>{\raggedright\arraybackslash}p{39.459mm}@{}}
\toprule
\textbf{Person and number} & \textbf{Present} & \textbf{Preterite} \\
\midrule
1 sg. & forceorfe & forcearf \\
2 sg. & forcierfst & forcurfe \\
3 sg. & forcierfð & forcearf \\
1–3 pl. & forceorfað & forcurfon \\
\bottomrule
\end{tabular}
\par\endgroup
% course-panel:T09-P01:end
% mdtable:T09:end
% mdparagraph:P019:begin
\textbf{Explanation:} The printed principal parts are \textbf{forceorfan — forcearf — forcurfon — forcorfen}. The selected present has \textbf{ceorf-} except for the second and third singular \textbf{cierf-} forms. The preterite first and third singular use \textbf{forcearf}; second singular \textbf{forcurfe} and all three plural persons use \textbf{curf-}, with \textbf{-e} and \textbf{-on} respectively. The past participle’s \textbf{corf-} is not the finite plural past stem. The selected \textbf{cierf} spelling is a declared rule-based normalization; attested \textbf{cirf/\allowbreak{}cyrf} variants do not justify arbitrary vowel substitutions. The principal parts are supplied diagnostics here, not a new test of independent nonfinite syntax. Cook 23.7’s printed reference \textbf{134} concerns \textbf{munan} in the selected grammar and is inapplicable to this strong III verb; the prompt remains unchanged. [3]
% mdparagraph:P019:end

\end{minipage}\par\addvspace{4pt}

\par\begin{minipage}{\linewidth}\raggedright
% mdparagraph:P020:begin
\textbf{F06. Commands and participle bases.}
% mdparagraph:P020:end

% item: NEW-U13-F06
% key: KEY-NEW-U13-F06
% mdtable:T10:begin
% course-table-style: 2026-10
% course-panel:T10-P01:begin
\par\begingroup
\fontsize{10.5}{13.2}\selectfont
\setlength{\LTleft}{0pt}\setlength{\LTright}{\fill}
\noindent
\begin{tabular}{@{}>{\raggedright\arraybackslash}p{30.902mm}>{\raggedright\arraybackslash}p{28.221mm}>{\raggedright\arraybackslash}p{28.221mm}>{\raggedright\arraybackslash}p{28.221mm}@{}}
\toprule
\textbf{Person and number} & \textbf{winnan} & \textbf{begān} & \textbf{behealdan} \\
\midrule
2 sg. & win & begā & beheald \\
2 pl. & winnað & begāð & behealdað \\
\bottomrule
\end{tabular}
\par\endgroup
% course-panel:T10-P01:end
% mdtable:T10:end
% mdtable:T11:begin
\par\begingroup
\fontsize{10.5}{13.2}\selectfont
\setlength{\LTleft}{0pt}\setlength{\LTright}{\fill}
\noindent
\begin{tabular}{@{}>{\raggedright\arraybackslash}p{32.04mm}>{\raggedright\arraybackslash}p{47.09mm}>{\raggedright\arraybackslash}p{39.25mm}@{}}
\toprule
\textbf{Verb} & \textbf{Present participle} & \textbf{Past participle} \\
\midrule
fordilgian & fordilgiende & fordilgod \\
onlīesan & onlīesende & onlīesed \\
\bottomrule
\end{tabular}
\par\endgroup

% mdtable:T11:end
% mdparagraph:P021:begin
\textbf{Explanation:} The singular commands address one person; \textbf{-að} or the contracted \textbf{-ð} pair supplies the plural. \textbf{Win} has a single final \textbf{n}, whereas \textbf{winnað} retains \textbf{nn} before the ending. \textbf{Beheald} keeps the unmutated present stem. Weak II \textbf{fordilgian} has \textbf{fordilgiende /\allowbreak{} fordilgod}; weak I \textbf{onlīesan} has \textbf{onlīesende /\allowbreak{} onlīesed}. These are participle bases, not full agreeing adjective tables. Do not substitute an indicative merely because its English translation could sound like a direction. [3]
% mdparagraph:P021:end

\end{minipage}\par\addvspace{4pt}

\reserve{7\baselineskip}
% mdparagraph:P022:begin
\section*{A. Clause analysis and reference}
% mdparagraph:P022:end

\reserve{5\baselineskip}
\par\begin{minipage}{\linewidth}\raggedright
% mdparagraph:P023:begin
\textbf{A01. Case under fronting.}
% mdparagraph:P023:end

% item: NEW-U13-A01
% key: KEY-NEW-U13-A01
% mdparagraph:P024:begin
Both clauses mean “The bishop greeted the monk.” In each, \textbf{se bisceop} is nominative subject, \textbf{ðone munuc} accusative direct object and \textbf{grētte} the finite preterite. The second clause fronts the complete object phrase without exchanging the participants. \textbf{Se/\allowbreak{}ðone} settle the distinct case roles even though both noun endings here are identical. Translating the second as “The monk greeted the bishop” confuses position with case.
% mdparagraph:P024:end

\end{minipage}\par\addvspace{4pt}

\par\begin{minipage}{\linewidth}\raggedright
% mdparagraph:P025:begin
\textbf{A02. Three functions of ðā.}
% mdparagraph:P025:end

% item: NEW-U13-A02
% key: KEY-NEW-U13-A02
% mdparagraph:P026:begin
“When the monk came, the bishop greeted the brothers.” The temporal clause is \textbf{[Ðā se munuc cōm]}: finite \textbf{cōm}, subject \textbf{se munuc}. The main clause is \textbf{[ðā grētte se bisceop ðā brōðor]}: finite \textbf{grētte}, subject \textbf{se bisceop}, object \textbf{ðā brōðor}. The first \textbf{ðā} is a temporal conjunction, the second its correlative adverb “then”, and the third the plural accusative article in the object phrase. Translating all three as “then” loses both the clause relation and the noun phrase.
% mdparagraph:P026:end

\end{minipage}\par\addvspace{4pt}

\par\begin{minipage}{\linewidth}\raggedright
% mdparagraph:P027:begin
\textbf{A03. A relative inside a main subject.}
% mdparagraph:P027:end

% item: NEW-U13-A03
% key: KEY-NEW-U13-A03
% mdparagraph:P028:begin
“The bishop whom the monk greeted sat in the house.” Main \textbf{sæt} has subject \textbf{se bisceop}. Relative \textbf{grētte} has subject \textbf{se munuc}. The antecedent \textbf{bisceop} is nominative as the main-clause subject but the person represented by \textbf{ðe} is the relative’s direct object. Invariant \textbf{ðe} does not display that accusative role by an ending. Do not turn the bishop into the greeter simply because his phrase comes first.
% mdparagraph:P028:end

\end{minipage}\par\addvspace{4pt}

\par\begin{minipage}{\linewidth}\raggedright
% mdparagraph:P029:begin
\textbf{A04. One continued subject.}
% mdparagraph:P029:end

% item: NEW-U13-A04
% key: KEY-NEW-U13-A04
% mdparagraph:P030:begin
“The monk greeted the bishop and gave him the food.” The monk performs both \textbf{grētte} and \textbf{sealde}; the second coordinated verb continues the first subject. \textbf{Ðone bisceop} is the accusative person greeted; \textbf{him} is the dative recipient, naturally the bishop; \textbf{ðone męte} is the accusative thing given. The sentence does not introduce the bishop as a new giver. As elsewhere, a larger context could clarify an unusual pronoun reference; the ordinary reading here follows the continued action.
% mdparagraph:P030:end

\end{minipage}\par\addvspace{4pt}

\par\begin{minipage}{\linewidth}\raggedright
% mdparagraph:P031:begin
\textbf{A05. A song and a person.}
% mdparagraph:P031:end

% item: NEW-U13-A05
% key: KEY-NEW-U13-A05
% mdparagraph:P032:begin
After \textbf{Se sang wæs swēte}, the second sentence means “It was pleasant”: \textbf{hē} refers to the song. After \textbf{Se bisceop cōm}, it means “He was in the house”: \textbf{hē} refers to the bishop. Masculine singular morphology permits either masculine noun in its respective context; lexical meaning and the actual preceding sentence identify the referent. English uses “it” for the song. Neither grammatical gender alone nor the spelling \textbf{hē} requires a human referent.
% mdparagraph:P032:end

\end{minipage}\par\addvspace{4pt}

\par\begin{minipage}{\linewidth}\raggedright
% mdparagraph:P033:begin
\textbf{A06. Real ambiguity.}
% mdparagraph:P033:end

% item: NEW-U13-A06
% key: KEY-NEW-U13-A06
% mdparagraph:P034:begin
Either the bishop or the monk may be the subject of \textbf{Hē sæt on ðǣm hūse}, “He sat in the house.” Both are masculine singular, so gender and number do not decide. Explicit alternatives are \textbf{Se bisceop sæt on ðǣm hūse.} and \textbf{Se munuc sæt on ðǣm hūse.} The task requires both rewrites. Choosing only the nearest noun, or calling one reading impossible without further context, overstates the evidence.
% mdparagraph:P034:end

\end{minipage}\par\addvspace{4pt}

\par\begin{minipage}{\linewidth}\raggedright
% mdparagraph:P035:begin
\textbf{A07. The quotation’s persons.}
% mdparagraph:P035:end

% item: NEW-U13-A07
% key: KEY-NEW-U13-A07
% mdparagraph:P036:begin
“The monk said, ‘I heard the song. Did you hear it?’” \textbf{Ic} is the monk, the quoted speaker. \textbf{Ðū} is the bishop, the singular addressee; \textbf{hīerdest} agrees with him. \textbf{Hine} is masculine accusative singular referring to \textbf{sang}, not to either man. First-singular \textbf{hīerde} agrees with \textbf{ic}. Changing the question to a plural address or translating \textbf{hine} as the monk changes the stated situation.
% mdparagraph:P036:end

\end{minipage}\par\addvspace{4pt}

\par\begin{minipage}{\linewidth}\raggedright
% mdparagraph:P037:begin
\textbf{A08. A report of the speaker’s own presence.}
% mdparagraph:P037:end

% item: NEW-U13-A08
% key: KEY-NEW-U13-A08
% mdparagraph:P038:begin
“The monk said that he was in the house.” The content clause begins at \textbf{ðæt} and continues through \textbf{wǣre}. Its finite verb is preterite subjunctive \textbf{wǣre}, with subject \textbf{hē}; the given context identifies that pronoun as the monk himself. The main verb is \textbf{sægde}, subject \textbf{se munuc}. The report does not assert that someone else was there, and the subjunctive does not by itself imply that the statement was false. “Had been” requires an appropriately prior reading of the context; no automatic English backshift rule decides it.
% mdparagraph:P038:end

\end{minipage}\par\addvspace{4pt}

\par\begin{minipage}{\linewidth}\raggedright
% mdparagraph:P039:begin
\textbf{A09. Restoring agreement.}
% mdparagraph:P039:end

% item: NEW-U13-A09
% key: KEY-NEW-U13-A09
% mdparagraph:P040:begin
“The bishop was in the church, and the brothers were in the house.” The full second clause is \textbf{Ðā brōðor wǣron on ðǣm hūse.} The omitted predicate is the corresponding past of “be”, restored with plural agreement for \textbf{ðā brōðor}. Copying singular \textbf{wæs} from the first clause would fail that agreement. A singular earlier predicate can support the understood meaning without supplying the literal plural form unchanged.
% mdparagraph:P040:end

\end{minipage}\par\addvspace{4pt}

\par\begin{minipage}{\linewidth}\raggedright
% mdparagraph:P041:begin
\textbf{A10. Initial phrases and clause boundaries.}
% mdparagraph:P041:end

% item: NEW-U13-A10
% key: KEY-NEW-U13-A10
% mdparagraph:P042:begin
First: “The monk sat in the house.” \textbf{On ðǣm hūse} is a fronted location phrase in the main clause; finite \textbf{sæt} has subject \textbf{se munuc}. Second: “When the monk sat in the house, the bishop came.” The initial temporal clause runs through \textbf{sæt}, with \textbf{se munuc} as subject; the main clause begins with correlative \textbf{ðā}, finite \textbf{cōm}, subject \textbf{se bisceop}. Its internal \textbf{on ðǣm hūse} belongs to the monk’s sitting, not the bishop’s coming. An initial phrase can be an adjunct or an entire dependent clause, so first position does not identify the actor.
% mdparagraph:P042:end

\end{minipage}\par\addvspace{4pt}

\par\begin{minipage}{\linewidth}\raggedright
% mdparagraph:P043:begin
\textbf{A11. A relative inside a temporal clause.}
% mdparagraph:P043:end

% item: NEW-U13-A11
% key: KEY-NEW-U13-A11
% mdparagraph:P044:begin
“When the monk whom the bishop greeted came, the brothers sang.” The structure is \textbf{[Temporal Ðā se munuc [Relative ðe se bisceop grētte] cōm], [Main ðā sungon ðā brōðor].} The subject of \textbf{grētte} is \textbf{se bisceop}; that of \textbf{cōm} is \textbf{se munuc}; that of \textbf{sungon} is \textbf{ðā brōðor}. Relative \textbf{ðe} refers to the monk and is the object of greeting. The first \textbf{ðā} is a conjunction, the second a correlative adverb, and the last a nominative plural article. The relative remains inside the temporal clause; it does not turn the bishop into the person who came.
% mdparagraph:P044:end

\end{minipage}\par\addvspace{4pt}

\par\begin{minipage}{\linewidth}\raggedright
% mdparagraph:P045:begin
\textbf{A12. Tendencies and explicit targets.}
% mdparagraph:P045:end

% item: NEW-U13-A12
% key: KEY-NEW-U13-A12
% mdparagraph:P046:begin
First, words and constituents differ. In \textbf{Ðone sang hīerde se munuc}, the two-word object phrase is one constituent, and the finite verb is the third written word. Moreover, the checked \textbf{Æfter ðon hē underðēodde} has a temporal phrase followed by a pronoun subject before the finite verb; even a second-constituent generalization is not compulsory. Second, dependent clauses commonly place their finite verb late, but the supplied exact \textbf{ðæt hē ǣrest gehīerde fram ēastsūðdǣle heofones} puts \textbf{gehīerde} before its source phrase. It means “that he first heard [the sound and song] from the south-eastern part of the sky”; the objects were fronted earlier in the full source sentence. This is a concrete dependent clause with a nonfinal finite verb. The rule describes a tendency, not a universal last-position law. Judge the complete clause and its dependencies, not a mechanical English-to-Old-English rearrangement. Cook 25.7 can still specify an exercise target: the final \textbf{bebēode} belongs to the immediately preceding relative. Other grammatical arrangements can be valid language while missing a specified target. The surrounding conditional’s \textbf{dōð} may precede the following relative or occupy a defensible final position; the nearest-relative interpretation is stated editorially. [4]
% mdparagraph:P046:end

\end{minipage}\par\addvspace{4pt}

\reserve{7\baselineskip}
% mdparagraph:P047:begin
\section*{C. Controlled composition}
% mdparagraph:P047:end

\reserve{5\baselineskip}
\par\begin{minipage}{\linewidth}\raggedright
% mdparagraph:P048:begin
\textbf{C01. Two initial constituents.}
% mdparagraph:P048:end

% item: NEW-U13-C01
% key: KEY-NEW-U13-C01
% mdparagraph:P049:begin
\itempara{1}{\textbf{Ðā brōðor hīerdon ðone sang.}}
% mdparagraph:P049:end

% mdparagraph:P050:begin
\itempara{2}{\textbf{Ðone sang hīerdon ðā brōðor.}}
% mdparagraph:P050:end

% mdparagraph:P051:begin
\textbf{Explanation:} Both models keep \textbf{ðā brōðor} nominative plural and \textbf{ðone sang} accusative singular. \textbf{Hīerdon} remains plural. The first begins with the subject, the second with the complete object phrase. Other valid orders do not replace the two requested arrangements; exchanging the article cases reverses the meaning.
% mdparagraph:P051:end

\end{minipage}\par\addvspace{4pt}

\par\begin{minipage}{\linewidth}\raggedright
% mdparagraph:P052:begin
\textbf{C02. A fronted temporal adverb.}
% mdparagraph:P052:end

% item: NEW-U13-C02
% key: KEY-NEW-U13-C02
% mdparagraph:P053:begin
\textbf{Ðā grētte se bisceop ðā brōðor.}
% mdparagraph:P053:end

% mdparagraph:P054:begin
\textbf{Explanation:} Initial \textbf{ðā} means “then”. The target places \textbf{grētte} before the nominal subject \textbf{se bisceop}; \textbf{ðā brōðor} is the accusative object. A subject-first sentence could be grammatical but would miss the expressly requested arrangement.
% mdparagraph:P054:end

\end{minipage}\par\addvspace{4pt}

\par\begin{minipage}{\linewidth}\raggedright
% mdparagraph:P055:begin
\textbf{C03. Temporal correlation.}
% mdparagraph:P055:end

% item: NEW-U13-C03
% key: KEY-NEW-U13-C03
% mdparagraph:P056:begin
\textbf{Ðā ðā brōðor cōmon, ðā sang se bisceop.}
% mdparagraph:P056:end

% mdparagraph:P057:begin
\textbf{Explanation:} The first \textbf{ðā} introduces “when”; \textbf{ðā brōðor} is its plural subject and \textbf{cōmon} its late finite verb. The next \textbf{ðā} is the main-clause correlative, followed by singular \textbf{sang} and its subject \textbf{se bisceop}. Do not give the brothers the singular past or make the bishop the subject of coming.
% mdparagraph:P057:end

\end{minipage}\par\addvspace{4pt}

\par\begin{minipage}{\linewidth}\raggedright
% mdparagraph:P058:begin
\textbf{C04. A relative object.}
% mdparagraph:P058:end

% item: NEW-U13-C04
% key: KEY-NEW-U13-C04
% mdparagraph:P059:begin
\textbf{Ðā brōðor, ðe se bisceop grētte, ēodon tō ðǣre cirican.}
% mdparagraph:P059:end

% mdparagraph:P060:begin
\textbf{Explanation:} The antecedent \textbf{ðā brōðor} is the main nominative plural subject of \textbf{ēodon}; inside the adjacent relative, \textbf{se bisceop} is the subject of singular \textbf{grētte}, and \textbf{ðe} represents its plural object. The unchanged particle does not copy the main subject’s case ending. A relative elsewhere may be intelligible, but this target requires it beside its antecedent.
% mdparagraph:P060:end

\end{minipage}\par\addvspace{4pt}

\par\begin{minipage}{\linewidth}\raggedright
% mdparagraph:P061:begin
\textbf{C05. Explicit reference.}
% mdparagraph:P061:end

% item: NEW-U13-C05
% key: KEY-NEW-U13-C05
% mdparagraph:P062:begin
\textbf{Se bisceop grētte ðone munuc. Ðā ēode se munuc tō ðǣre cirican.}
% mdparagraph:P062:end

% mdparagraph:P063:begin
\textbf{Explanation:} The bishop performs the greeting; the monk performs the later movement. Repeating \textbf{se munuc} makes the second subject explicit as requested. A bare \textbf{hē} might be grammatical in a fuller context but would fail this clarity target. \textbf{Ðā} here means “then”, not “when”.
% mdparagraph:P063:end

\end{minipage}\par\addvspace{4pt}

\par\begin{minipage}{\linewidth}\raggedright
% mdparagraph:P064:begin
\textbf{C06. A singular speaker and plural addressees.}
% mdparagraph:P064:end

% item: NEW-U13-C06
% key: KEY-NEW-U13-C06
% mdparagraph:P065:begin
\textbf{Se munuc cwæð: “Ic wille mid ēow sprecan.”}
% mdparagraph:P065:end

% mdparagraph:P066:begin
\textbf{Explanation:} \textbf{Ic} refers to the monk and takes \textbf{wille}; \textbf{ēow} is plural because he addresses several brothers. \textbf{Mid ēow sprecan} uses the taught conversation frame and plain infinitive after \textbf{willan}. Singular \textbf{ðē} would change the specified audience. Keep the words inside the quotation in the speaker’s own person.
% mdparagraph:P066:end

\end{minipage}\par\addvspace{4pt}

\par\begin{minipage}{\linewidth}\raggedright
% mdparagraph:P067:begin
\textbf{C07. Coordination with a continued subject.}
% mdparagraph:P067:end

% item: NEW-U13-C07
% key: KEY-NEW-U13-C07
% mdparagraph:P068:begin
\textbf{Se munuc geseah ðone bisceop and grētte hine.}
% mdparagraph:P068:end

% mdparagraph:P069:begin
\textbf{Explanation:} The monk sees and greets; \textbf{ðone bisceop} and \textbf{hine} are accusative objects referring to the bishop. The coordinated \textbf{grētte} continues the single expressed subject \textbf{se munuc}. Repeating that subject can be grammatical, but this exercise explicitly asks for one expression of it.
% mdparagraph:P069:end

\end{minipage}\par\addvspace{4pt}

\par\begin{minipage}{\linewidth}\raggedright
% mdparagraph:P070:begin
\textbf{C08. Three nested clauses.}
% mdparagraph:P070:end

% item: NEW-U13-C08
% key: KEY-NEW-U13-C08
% mdparagraph:P071:begin
\textbf{Ðā se bisceop, ðe se munuc grētte, cōm, ðā sungon ðā brōðor.}
% mdparagraph:P071:end

% mdparagraph:P072:begin
\textbf{Explanation:} The clause map is \textbf{[Temporal Ðā se bisceop [Relative ðe se munuc grētte] cōm], [Main ðā sungon ðā brōðor].} The monk is subject of \textbf{grētte}, the bishop of \textbf{cōm}, and the brothers of \textbf{sungon}. Relative \textbf{ðe} is the object of greeting and refers to the bishop. Preserve the temporal/\allowbreak{}main correlation and mark all three clauses; identifying only the relative and the main leaves one requested analysis missing.
% mdparagraph:P072:end

\end{minipage}\par\addvspace{4pt}

\reserve{7\baselineskip}
% mdparagraph:P073:begin
\section*{Cook Group VII. Exercise 22}
% mdparagraph:P073:end

\reserve{5\baselineskip}
\par\begin{minipage}{\linewidth}\raggedright
% mdparagraph:P074:begin
\textbf{22.1.}
% mdparagraph:P074:end

% item: EX-VII-22-01
% key: KEY-EX-VII-22-01
% mdparagraph:P075:begin
\textbf{Prompt:} Idleness is the mother of sin.
% mdparagraph:P075:end

% mdparagraph:P076:begin
\textbf{Model:} Īdelnes is synne mōdor.
% mdparagraph:P076:end

% mdparagraph:P077:begin
\textbf{Explanation:} \textbf{Īdelnes} is the subject; \textbf{mōdor} is the predicate nominative. Genitive \textbf{synne} depends on “mother”. Both nouns express general ideas without a compulsory article. The selected nominative \textbf{īdelnes} has one final \textbf{s}, while its oblique stem has \textbf{ss}. Making “sin” the subject reverses the relation.
% mdparagraph:P077:end

\end{minipage}\par\addvspace{4pt}

\par\begin{minipage}{\linewidth}\raggedright
% mdparagraph:P078:begin
\textbf{22.2.}
% mdparagraph:P078:end

% item: EX-VII-22-02
% key: KEY-EX-VII-22-02
% mdparagraph:P079:begin
\textbf{Prompt:} Etheldred was (om. the) daughter of the king of the East Anglians.
% mdparagraph:P079:end

% mdparagraph:P080:begin
\textbf{Model:} Æðelðrȳð wæs Eastęngla cyninges dohtor.
% mdparagraph:P080:end

% mdparagraph:P081:begin
\textbf{Explanation:} The supplied nominative name is \textbf{Æðelðrȳð}. The nested genitives mean “daughter [of the king [of the East Anglians]]”: \textbf{Eastęngla} depends on \textbf{cyninges}, and that king phrase depends on predicate \textbf{dohtor}. Cook expressly omits the article before “daughter”. Preserve the selected short \textbf{o} of \textbf{dohtor}; no inflection of the supplied name is tested.
% mdparagraph:P081:end

\end{minipage}\par\addvspace{4pt}

\par\begin{minipage}{\linewidth}\raggedright
% mdparagraph:P082:begin
\textbf{22.3.}
% mdparagraph:P082:end

% item: EX-VII-22-03
% key: KEY-EX-VII-22-03
% mdparagraph:P083:begin
\textbf{Prompt:} The bishop’s beautiful monastery was destroyed (fordilgian).
% mdparagraph:P083:end

% mdparagraph:P084:begin
\textbf{Model:} Ðæt fægere mynster ðæs bisceopes wæs fordilgod.
% mdparagraph:P084:end

% mdparagraph:P085:begin
\textbf{Explanation:} The monastery is the patient of \textbf{wæs fordilgod}, a preterite passive, not the agent of destruction. \textbf{Mynster} is neuter singular; the overt article \textbf{ðæt} licenses weak \textbf{fægere}. The following genitive \textbf{ðæs bisceopes} identifies the possessor. This uses the already taught article-and-adjective pattern without assuming a general rule after limiting noun genitives. \textbf{Fordilgod} is the separately taught past-participle base. Do not translate with an active “destroyed” unless the monastery is wrongly made the destroyer.
% mdparagraph:P085:end

\end{minipage}\par\addvspace{4pt}

\par\begin{minipage}{\linewidth}\raggedright
% mdparagraph:P086:begin
\textbf{22.4.}
% mdparagraph:P086:end

% item: EX-VII-22-04
% key: KEY-EX-VII-22-04
% mdparagraph:P087:begin
\textbf{Prompt:} Labor [toil, sing.] with your hands.
% mdparagraph:P087:end

% mdparagraph:P088:begin
\textbf{Model:} Win mid ðīnum handum.
% mdparagraph:P088:end

% mdparagraph:P089:begin
\textbf{Explanation:} Cook specifies a singular addressee: \textbf{win}. Possessive \textbf{ðīnum} is therefore singular “your”, though it agrees with plural dative \textbf{handum} after \textbf{mid}. \textbf{Winnað} would address several people and miss this control. The command’s singular \textbf{n} differs from the plural stem’s \textbf{nn}.
% mdparagraph:P089:end

\end{minipage}\par\addvspace{4pt}

\par\begin{minipage}{\linewidth}\raggedright
% mdparagraph:P090:begin
\textbf{22.5.}
% mdparagraph:P090:end

% item: EX-VII-22-05
% key: KEY-EX-VII-22-05
% mdparagraph:P091:begin
\textbf{Prompt:} Pursue (sing.) your study.
% mdparagraph:P091:end

% mdparagraph:P092:begin
\textbf{Model:} Begā ðīne leornunge.
% mdparagraph:P092:end

% mdparagraph:P093:begin
\textbf{Explanation:} \textbf{Begā} is the singular command “pursue”; \textbf{ðīne leornunge} is its feminine accusative object. The final \textbf{-e} belongs to agreement and the noun’s oblique system, not to a plural addressee. A form meaning simply “go” would miss the lexical use specified here.
% mdparagraph:P093:end

\end{minipage}\par\addvspace{4pt}

\par\begin{minipage}{\linewidth}\raggedright
% mdparagraph:P094:begin
\textbf{22.6.}
% mdparagraph:P094:end

% item: EX-VII-22-06
% key: KEY-EX-VII-22-06
% mdparagraph:P095:begin
\textbf{Prompt:} He goes to church.
% mdparagraph:P095:end

% mdparagraph:P096:begin
\textbf{Model:} Hē gǣð tō cirican.
% mdparagraph:P096:end

% mdparagraph:P097:begin
\textbf{Explanation:} \textbf{Gǣð} is present third singular, agreeing with \textbf{hē}. \textbf{Tō} takes the dative \textbf{cirican} in the selected ordinary weak feminine system. The customary church expression need not add an article. The exact reading’s \textbf{ciricean} is a separately supplied spelling, not a reason to alter the selected productive table.
% mdparagraph:P097:end

\end{minipage}\par\addvspace{4pt}

\par\begin{minipage}{\linewidth}\raggedright
% mdparagraph:P098:begin
\textbf{22.7.}
% mdparagraph:P098:end

% item: EX-VII-22-07
% key: KEY-EX-VII-22-07
% mdparagraph:P099:begin
\textbf{Prompt:} Depart from me, ye that work unrighteousness.
% mdparagraph:P099:end

% mdparagraph:P100:begin
\textbf{Model:} Gewītað fram mē, gē ðe wyrceað unrihtwīsnesse.
% mdparagraph:P100:end

% mdparagraph:P101:begin
\textbf{Explanation:} \textbf{Gewītað} addresses the plural \textbf{gē}; \textbf{fram mē} has a dative complement. The relative \textbf{ðe wyrceað} keeps that plural reference, and \textbf{unrihtwīsnesse} is the accusative thing worked or done. The already taught ordinary sense of \textbf{wyrcean} covers doing a deed; this is not the distinct new “make four parts” use in 23.4. Singular commands or finite endings would change the audience.
% mdparagraph:P101:end

\end{minipage}\par\addvspace{4pt}

\par\begin{minipage}{\linewidth}\raggedright
% mdparagraph:P102:begin
\textbf{22.8.}
% mdparagraph:P102:end

% item: EX-VII-22-08
% key: KEY-EX-VII-22-08
% mdparagraph:P103:begin
\textbf{Prompt:} What is your name?
% mdparagraph:P103:end

% mdparagraph:P104:begin
\textbf{Model:} Hwæt is ðīn nama?
% mdparagraph:P104:end

% mdparagraph:P105:begin
\textbf{Explanation:} \textbf{Hwæt} asks the predicate identity; \textbf{ðīn nama} is the masculine nominative subject of \textbf{is}. \textbf{Nama} remains nominative in a question. The model selects singular “your”; if no external context specifies the audience, a correctly formed plural address is also defensible. English interrogative order does not supply an Old English accusative for “name”.
% mdparagraph:P105:end

\end{minipage}\par\addvspace{4pt}

\par\begin{minipage}{\linewidth}\raggedright
% mdparagraph:P106:begin
\textbf{22.9.}
% mdparagraph:P106:end

% item: EX-VII-22-09
% key: KEY-EX-VII-22-09
% mdparagraph:P107:begin
\textbf{Prompt:} Tell (hātan) our three friends to go thither to him.
% mdparagraph:P107:end

% mdparagraph:P108:begin
\textbf{Model:} Hāt ūre ðrīe frīend ðider tō him gān.
% mdparagraph:P108:end

% mdparagraph:P109:begin
\textbf{Explanation:} \textbf{Hāt} introduces the taught accusative-person plus infinitive construction. \textbf{Ūre ðrīe frīend} is the plural accusative actor of \textbf{gān}; \textbf{ðider} gives “thither”, and \textbf{tō him} names the destination person. The model addresses one person; \textbf{hātað} is a valid plural command if several people are being addressed, since this prompt does not specify their number. Do not make the three friends the subject of the command itself.
% mdparagraph:P109:end

\end{minipage}\par\addvspace{4pt}

\par\begin{minipage}{\linewidth}\raggedright
% mdparagraph:P110:begin
\textbf{22.10.}
% mdparagraph:P110:end

% item: EX-VII-22-10
% key: KEY-EX-VII-22-10
% mdparagraph:P111:begin
\textbf{Prompt:} The song descends from heaven (plur.) to earth.
% mdparagraph:P111:end

% mdparagraph:P112:begin
\textbf{Model:} Se sang āstīhð of heofonum tō eorðan.
% mdparagraph:P112:end

% mdparagraph:P113:begin
\textbf{Explanation:} The singular song takes \textbf{āstīhð}. The chapter’s selected sense permits “descend”; movement downward is fixed by \textbf{of heofonum tō eorðan}, not by changing the verb to an invented form. Cook requests plural heaven, hence dative plural \textbf{heofonum} after \textbf{of}. \textbf{Eorðan} is dative after \textbf{tō}.
% mdparagraph:P113:end

\end{minipage}\par\addvspace{4pt}

\par\begin{minipage}{\linewidth}\raggedright
% mdparagraph:P114:begin
\textbf{22.11.}
% mdparagraph:P114:end

% item: EX-VII-22-11
% key: KEY-EX-VII-22-11
% mdparagraph:P115:begin
\textbf{Prompt:} I am wont to visit the bishop.
% mdparagraph:P115:end

% mdparagraph:P116:begin
\textbf{Model:} Ic gewunie ðone bisceop genēosian.
% mdparagraph:P116:end

% mdparagraph:P117:begin
\textbf{Explanation:} First-singular \textbf{gewunie} governs the plain infinitive \textbf{genēosian} in the taught habitual frame, “am accustomed to”. The bishop is its accusative object. The present expresses a habit; it does not need a separate copula corresponding word for word to English “am wont”. A past \textbf{gewunode} would change the tense.
% mdparagraph:P117:end

\end{minipage}\par\addvspace{4pt}

\par\begin{minipage}{\linewidth}\raggedright
% mdparagraph:P118:begin
\textbf{22.12.}
% mdparagraph:P118:end

% item: EX-VII-22-12
% key: KEY-EX-VII-22-12
% mdparagraph:P119:begin
\textbf{Prompt:} I wish (willan) first (om. to) ask you (sing.) one thing (156).
% mdparagraph:P119:end

% mdparagraph:P120:begin
\textbf{Model:} Ic wille ǣrest ðē āscian ānes ðinges.
% mdparagraph:P120:end

% mdparagraph:P121:begin
\textbf{Explanation:} \textbf{Wille} is first singular and takes plain \textbf{āscian}, without \textbf{tō}. \textbf{Ǣrest} means “first”. In the selected asking frame the person \textbf{ðē} is accusative and the thing requested is genitive \textbf{ānes ðinges}. The spelling \textbf{ðē} alone could fit more than one case, but the frame fixes its function here. An accusative “one thing” misses Cook’s specified genitive target.
% mdparagraph:P121:end

\end{minipage}\par\addvspace{4pt}

\par\begin{minipage}{\linewidth}\raggedright
% mdparagraph:P122:begin
\textbf{22.13.}
% mdparagraph:P122:end

% item: EX-VII-22-13
% key: KEY-EX-VII-22-13
% mdparagraph:P123:begin
\textbf{Prompt:} I wish to take (lǣdan) him with me.
% mdparagraph:P123:end

% mdparagraph:P124:begin
\textbf{Model:} Ic wille hine mid mē lǣdan.
% mdparagraph:P124:end

% mdparagraph:P125:begin
\textbf{Explanation:} \textbf{Wille} takes the newly assigned plain infinitive \textbf{lǣdan}. Accusative \textbf{hine} is the person taken or led; dative \textbf{mē} belongs to \textbf{mid}, “with me”. The action is taking someone along, not merely receiving or seizing him. Substituting \textbf{niman} would miss Cook’s lexical control even where another context could make it grammatical.
% mdparagraph:P125:end

\end{minipage}\par\addvspace{4pt}

\par\begin{minipage}{\linewidth}\raggedright
% mdparagraph:P126:begin
\textbf{22.14.}
% mdparagraph:P126:end

% item: EX-VII-22-14
% key: KEY-EX-VII-22-14
% mdparagraph:P127:begin
\textbf{Prompt:} The priest prostrated himself before (beforan, 172) the bishop.
% mdparagraph:P127:end

% mdparagraph:P128:begin
\textbf{Model:} Se mæsseprēost āstreahte hine beforan ðone bisceop.
% mdparagraph:P128:end

% mdparagraph:P129:begin
\textbf{Explanation:} The priest is subject of preterite \textbf{āstreahte}; accusative \textbf{hine} is reflexive through reference to that same subject. The model construes \textbf{beforan ðone bisceop} as positioning before the bishop, supported by the checked accusative use. A static \textbf{beforan ðǣm bisceope} is a defensible alternative construal. Do not claim that every occurrence of \textbf{beforan} requires motion or accusative. [5]
% mdparagraph:P129:end

\end{minipage}\par\addvspace{4pt}

\par\begin{minipage}{\linewidth}\raggedright
% mdparagraph:P130:begin
\textbf{22.15.}
% mdparagraph:P130:end

% item: EX-VII-22-15
% key: KEY-EX-VII-22-15
% mdparagraph:P131:begin
\textbf{Prompt:} On the third day the holy soul was released from the body.
% mdparagraph:P131:end

% mdparagraph:P132:begin
\textbf{Model:} On ðone ðriddan dæg wæs sēo hālge sāwol of ðǣm līchaman onlīesed.
% mdparagraph:P132:end

% mdparagraph:P133:begin
\textbf{Explanation:} \textbf{On ðone ðriddan dæg} is the specified time phrase; the ordinal has weak \textbf{-an} after the article. \textbf{Sēo hālge sāwol} is the feminine patient of \textbf{wæs … onlīesed}. The selected weak adjective has the audited syncopated \textbf{hālg-} stem. \textbf{Of ðǣm līchaman} means “from the body”, with dative after \textbf{of}. An active finite release verb would require a different subject relation.
% mdparagraph:P133:end

\end{minipage}\par\addvspace{4pt}

\reserve{7\baselineskip}
% mdparagraph:P134:begin
\section*{Cook Group VII. Exercise 23}
% mdparagraph:P134:end

\reserve{5\baselineskip}
\par\begin{minipage}{\linewidth}\raggedright
% mdparagraph:P135:begin
\textbf{23.1.}
% mdparagraph:P135:end

% item: EX-VII-23-01
% key: KEY-EX-VII-23-01
% mdparagraph:P136:begin
\textbf{Prompt:} The largest window of the church is two feet (154. c) broad, and four feet high.
% mdparagraph:P136:end

% mdparagraph:P137:begin
\textbf{Model:} Ðæt mǣste ēagðȳrel ðǣre cirican is twēgra fōta brād and fēower fōta hēah.
% mdparagraph:P137:end

% mdparagraph:P138:begin
\textbf{Explanation:} The subject is neuter \textbf{ðæt mǣste ēagðȳrel}. Genitive \textbf{ðǣre cirican} identifies its church. The measures are genitive \textbf{twēgra fōta} and \textbf{fēower fōta}, dependent on \textbf{brād} and \textbf{hēah} respectively. Both adjectives describe the window. Ordinary adjective order can vary; nominative or accusative feet would miss the assigned measure construction.
% mdparagraph:P138:end

\end{minipage}\par\addvspace{4pt}

\par\begin{minipage}{\linewidth}\raggedright
% mdparagraph:P139:begin
\textbf{23.2.}
% mdparagraph:P139:end

% item: EX-VII-23-02
% key: KEY-EX-VII-23-02
% mdparagraph:P140:begin
\textbf{Prompt:} Little by little the sweet song drew nigh.
% mdparagraph:P140:end

% mdparagraph:P141:begin
\textbf{Model:} Styccemǣlum nēalǣcte se swēta sang.
% mdparagraph:P141:end

% mdparagraph:P142:begin
\textbf{Explanation:} Initial \textbf{styccemǣlum} means “little by little”; the finite singular past is \textbf{nēalǣcte}, with \textbf{ct}, not an assumed \textbf{ht}. \textbf{Se swēta sang} has weak nominative masculine \textbf{swēta} after its article. The fronted adverb does not become a subject, and the pleasant sound’s masculine gender does not make it a person.
% mdparagraph:P142:end

\end{minipage}\par\addvspace{4pt}

\par\begin{minipage}{\linewidth}\raggedright
% mdparagraph:P143:begin
\textbf{23.3.}
% mdparagraph:P143:end

% item: EX-VII-23-03
% key: KEY-EX-VII-23-03
% mdparagraph:P144:begin
\textbf{Prompt:} I entreat you to tell [that you tell] me what is (194. b) the angel’s name.
% mdparagraph:P144:end

% mdparagraph:P145:begin
\textbf{Model:} Ic hālsie ðē ðæt ðū mē sęcge hwæt ðæs ęngles nama sīe.
% mdparagraph:P145:end

% mdparagraph:P146:begin
\textbf{Explanation:} \textbf{Hālsie} takes accusative \textbf{ðē} as the person entreated. The request content is \textbf{ðæt ðū mē sęcge}: present subjunctive \textbf{sęcge}, dative recipient \textbf{mē}. Inside it, \textbf{hwæt ðæs ęngles nama sīe} is the indirect question, with subjunctive \textbf{sīe}. There are two dependent clauses and two different subjects to trace. A pair of detached direct questions would miss the requested construction.
% mdparagraph:P146:end

\end{minipage}\par\addvspace{4pt}

\par\begin{minipage}{\linewidth}\raggedright
% mdparagraph:P147:begin
\textbf{23.4.}
% mdparagraph:P147:end

% item: EX-VII-23-04
% key: KEY-EX-VII-23-04
% mdparagraph:P148:begin
\textbf{Prompt:} The soldiers took his raiment, and made (wyrcean) four parts, (om. to) every soldier a part; thus did the soldiers.
% mdparagraph:P148:end

% mdparagraph:P149:begin
\textbf{Model:} Ðā cęmpan nōmon his rēaf and worhton fēower dǣlas, ǣlcum cęmpan ānne dǣl; swā dydon ðā cęmpan.
% mdparagraph:P149:end

% mdparagraph:P150:begin
\textbf{Explanation:} Both \textbf{nōmon} and \textbf{worhton} have plural subject \textbf{ðā cęmpan}. \textbf{His rēaf} is what they took; \textbf{fēower dǣlas} is what they made. \textbf{Ǣlcum cęmpan} is dative recipient, while \textbf{ānne dǣl} is the part allotted, without an added preposition for English “to”. The final \textbf{swā dydon ðā cęmpan} resumes the soldiers explicitly. Singular \textbf{worhte} would break agreement.
% mdparagraph:P150:end

\end{minipage}\par\addvspace{4pt}

\par\begin{minipage}{\linewidth}\raggedright
% mdparagraph:P151:begin
\textbf{23.5.}
% mdparagraph:P151:end

% item: EX-VII-23-05
% key: KEY-EX-VII-23-05
% mdparagraph:P152:begin
\textbf{Prompt:} The mother of Jesus stood by (wið, 172) the cross.
% mdparagraph:P152:end

% mdparagraph:P153:begin
\textbf{Model:} Ðæs Hǣlendes mōdor stōd wið ðā rōde.
% mdparagraph:P153:end

% mdparagraph:P154:begin
\textbf{Explanation:} \textbf{Ðæs Hǣlendes} depends on nominative \textbf{mōdor}; she is the subject of \textbf{stōd}. The model’s \textbf{wið ðā rōde} uses the independently attested static proximity frame with accusative. It means standing beside the cross, not moving towards it. A defensible attested dative construal is not inherently bad language, but the selected model illustrates the assigned accusative frame. [5]
% mdparagraph:P154:end

\end{minipage}\par\addvspace{4pt}

\par\begin{minipage}{\linewidth}\raggedright
% mdparagraph:P155:begin
\textbf{23.6.}
% mdparagraph:P155:end

% item: EX-VII-23-06
% key: KEY-EX-VII-23-06
% mdparagraph:P156:begin
\textbf{Prompt:} Many read this writing.
% mdparagraph:P156:end

% mdparagraph:P157:begin
\textbf{Model:} Manige rǣdað ðis gewrit.
% mdparagraph:P157:end

% mdparagraph:P158:begin
\textbf{Explanation:} Substantival plural \textbf{manige}, “many people”, takes plural \textbf{rǣdað}. \textbf{Ðis gewrit} is neuter singular accusative, “this writing”. The number of readers does not make the writing plural. Written English “read” also permits a past interpretation: \textbf{Manige rǣddon ðis gewrit.} A singular “many a person” construal would require corresponding agreement and a different subject pattern; it is not the selected plural construction.
% mdparagraph:P158:end

\end{minipage}\par\addvspace{4pt}

\par\begin{minipage}{\linewidth}\raggedright
% mdparagraph:P159:begin
\textbf{23.7.}
% mdparagraph:P159:end

% item: EX-VII-23-07
% key: KEY-EX-VII-23-07
% mdparagraph:P160:begin
\textbf{Prompt:} With his little hatchet he cut down (forceorfan, 134) the tree.
% mdparagraph:P160:end

% mdparagraph:P161:begin
\textbf{Model:} Mid his lȳtlan adesan forcearf hē ðæt trēo.
% mdparagraph:P161:end

% mdparagraph:P162:begin
\textbf{Explanation:} \textbf{Mid his lȳtlan adesan} is a dative instrumental phrase; weak \textbf{lȳtlan} agrees with masculine \textbf{adesan}. Singular \textbf{forcearf} has subject \textbf{hē}, and \textbf{ðæt trēo} is its accusative object. Cook’s numeral \textbf{134} is retained, with its inapplicability documented; this is the audited strong III past. Reading the initial hatchet phrase as subject or using plural \textbf{forcurfon} changes the construction. [3]
% mdparagraph:P162:end

\end{minipage}\par\addvspace{4pt}

\par\begin{minipage}{\linewidth}\raggedright
% mdparagraph:P163:begin
\textbf{23.8.}
% mdparagraph:P163:end

% item: EX-VII-23-08
% key: KEY-EX-VII-23-08
% mdparagraph:P164:begin
\textbf{Prompt:} The ax is placed (āsęttan) into (tō) the root (wyrttruma, 53) of the trees.
% mdparagraph:P164:end

% mdparagraph:P165:begin
\textbf{Model:} Sēo æx is tō ðǣra trēowa wyrttruman āsętt.
% mdparagraph:P165:end

% mdparagraph:P166:begin
\textbf{Explanation:} Feminine \textbf{sēo æx} is the patient of present passive \textbf{is … āsętt}. The locally supplied placement phrase uses \textbf{tō} with dative \textbf{wyrttruman}, modified by genitive plural \textbf{ðǣra trēowa}. Ordinary \textbf{ðāra} is a valid genitive plural variant. Keep this bounded at/\allowbreak{}to-root expression; it does not establish that \textbf{tō} generally marks static location.
% mdparagraph:P166:end

\end{minipage}\par\addvspace{4pt}

\par\begin{minipage}{\linewidth}\raggedright
% mdparagraph:P167:begin
\textbf{23.9.}
% mdparagraph:P167:end

% item: EX-VII-23-09
% key: KEY-EX-VII-23-09
% mdparagraph:P168:begin
\textbf{Prompt:} He took an ax in his hand.
% mdparagraph:P168:end

% mdparagraph:P169:begin
\textbf{Model:} Hē nōm æxe on his hand.
% mdparagraph:P169:end

% mdparagraph:P170:begin
\textbf{Explanation:} \textbf{Nōm} is singular preterite; accusative \textbf{æxe} is the axe taken. The supplied \textbf{on his hand} construes the result as taking it “into his hand”, with accusative \textbf{hand}. A static dative \textbf{on his handa}, “in his hand”, is defensible where that is the intended construal. This model is declared teaching composition, jointly supported by the qualified directional rule and the lexical sense of \textbf{niman}; it is not a verbatim source quotation. [5]
% mdparagraph:P170:end

\end{minipage}\par\addvspace{4pt}

\par\begin{minipage}{\linewidth}\raggedright
% mdparagraph:P171:begin
\textbf{23.10.}
% mdparagraph:P171:end

% item: EX-VII-23-10
% key: KEY-EX-VII-23-10
% mdparagraph:P172:begin
\textbf{Prompt:} Sell (sęllan) me five hatchets.
% mdparagraph:P172:end

% mdparagraph:P173:begin
\textbf{Model:} Sęle mē fīf adesan.
% mdparagraph:P173:end

% mdparagraph:P174:begin
\textbf{Explanation:} Singular command \textbf{sęle} takes dative \textbf{mē} as buyer/\allowbreak{}recipient and accusative plural \textbf{fīf adesan} as what is sold. The lexical sense here is “sell”, not merely “give”. A plural command \textbf{sęllað} is possible if several sellers are addressed, since the English prompt leaves the addressee’s number open. Do not give the five hatchets a singular noun form.
% mdparagraph:P174:end

\end{minipage}\par\addvspace{4pt}

\par\begin{minipage}{\linewidth}\raggedright
% mdparagraph:P175:begin
\textbf{23.11.}
% mdparagraph:P175:end

% item: EX-VII-23-11
% key: KEY-EX-VII-23-11
% mdparagraph:P176:begin
\textbf{Prompt:} His bodily infirmity was very grievous.
% mdparagraph:P176:end

% mdparagraph:P177:begin
\textbf{Model:} His līchamlicu untrymnes wæs swīðe hęfigu.
% mdparagraph:P177:end

% mdparagraph:P178:begin
\textbf{Explanation:} \textbf{Untrymnes} is feminine singular. Attributive \textbf{līchamlicu} and predicative \textbf{hęfigu} are the selected strong feminine nominative forms; \textbf{his} identifies whose infirmity it is. \textbf{Swīðe} is the adverb “very”. A masculine or neuter adjective ending would not agree with the actual noun. The past copula \textbf{wæs} preserves the prompt’s time.
% mdparagraph:P178:end

\end{minipage}\par\addvspace{4pt}

\par\begin{minipage}{\linewidth}\raggedright
% mdparagraph:P179:begin
\textbf{23.12.}
% mdparagraph:P179:end

% item: EX-VII-23-12
% key: KEY-EX-VII-23-12
% mdparagraph:P180:begin
\textbf{Prompt:} Your companionship is very pleasant.
% mdparagraph:P180:end

% mdparagraph:P181:begin
\textbf{Model:} Ðīn gefērscipe is swīðe wynsum.
% mdparagraph:P181:end

% mdparagraph:P182:begin
\textbf{Explanation:} Masculine \textbf{gefērscipe} is the subject and \textbf{wynsum} its predicate adjective; \textbf{swīðe} supplies “very”. The model takes “your” as singular \textbf{ðīn}. Plural possessive \textbf{ēower} is defensible for a plural addressee if context requires it. Do not transfer a feminine \textbf{-u} from the preceding exercise to this masculine noun.
% mdparagraph:P182:end

\end{minipage}\par\addvspace{4pt}

\par\begin{minipage}{\linewidth}\raggedright
% mdparagraph:P183:begin
\textbf{23.13.}
% mdparagraph:P183:end

% item: EX-VII-23-13
% key: KEY-EX-VII-23-13
% mdparagraph:P184:begin
\textbf{Prompt:} Is it any wonder that [though] the faith of God grew (196. e)?
% mdparagraph:P184:end

% mdparagraph:P185:begin
\textbf{Model:} Is hit ǣnig wundor ðēah Godes gelēafa wēoxe?
% mdparagraph:P185:end

% mdparagraph:P186:begin
\textbf{Explanation:} The supplied wonder construction has anticipatory \textbf{hit}, predicate \textbf{ǣnig wundor}, and a \textbf{ðēah} clause giving the conceded circumstance. \textbf{Godes gelēafa} is the subject of preterite subjunctive \textbf{wēoxe}. An admitted event may occur in this concessive subjunctive; translating it as necessarily unreal is unwarranted. The literal “though” control prevents replacing the structure merely with an unexamined English “that”.
% mdparagraph:P186:end

\end{minipage}\par\addvspace{4pt}

\reserve{7\baselineskip}
% mdparagraph:P187:begin
\section*{Cook Group VII. Exercise 24}
% mdparagraph:P187:end

\reserve{5\baselineskip}
\par\begin{minipage}{\linewidth}\raggedright
% mdparagraph:P188:begin
\textbf{24.1.}
% mdparagraph:P188:end

% item: EX-VII-24-01
% key: KEY-EX-VII-24-01
% mdparagraph:P189:begin
\textbf{Prompt:} He that (sē ðe) sitteth on the throne shall dwell (wunian) in the midst of them.
% mdparagraph:P189:end

% mdparagraph:P190:begin
\textbf{Model:} Sē ðe on ðǣm setle sitt, hē wunað tōmiddes him.
% mdparagraph:P190:end

% mdparagraph:P191:begin
\textbf{Explanation:} The initial \textbf{sē ðe} clause identifies the sitter; \textbf{on ðǣm setle} is static dative. Resumptive \textbf{hē} begins the main predication, with present \textbf{wunað} understood prospectively here. \textbf{Tōmiddes him} means “in their midst”; the pronoun’s dative plural role follows its antecedent “them”. The future English does not compel an extra auxiliary. Ordinary short \textbf{se} is a taught spelling alternative to the supplied \textbf{sē}.
% mdparagraph:P191:end

\end{minipage}\par\addvspace{4pt}

\par\begin{minipage}{\linewidth}\raggedright
% mdparagraph:P192:begin
\textbf{24.2.}
% mdparagraph:P192:end

% item: EX-VII-24-02
% key: KEY-EX-VII-24-02
% mdparagraph:P193:begin
\textbf{Prompt:} Angels are spirits without body.
% mdparagraph:P193:end

% mdparagraph:P194:begin
\textbf{Model:} Ęnglas sindon gāstas būtan līchaman.
% mdparagraph:P194:end

% mdparagraph:P195:begin
\textbf{Explanation:} Both \textbf{ęnglas} and predicate \textbf{gāstas} are nominative plural, linked by plural \textbf{sindon}. \textbf{Būtan līchaman} uses the taught dative “without body” frame. A singular copula would fail agreement; making “spirits” accusative would misanalyse the predicate. Another taught plural form of “be” can be valid if its ordinary use fits the statement.
% mdparagraph:P195:end

\end{minipage}\par\addvspace{4pt}

\par\begin{minipage}{\linewidth}\raggedright
% mdparagraph:P196:begin
\textbf{24.3.}
% mdparagraph:P196:end

% item: EX-VII-24-03
% key: KEY-EX-VII-24-03
% mdparagraph:P197:begin
\textbf{Prompt:} I saw a star (169) fall from heaven unto [until] the earth.
% mdparagraph:P197:end

% mdparagraph:P198:begin
\textbf{Model:} Ic geseah steorran of heofone oð eorðan feallan.
% mdparagraph:P198:end

% mdparagraph:P199:begin
\textbf{Explanation:} \textbf{Geseah} is first-singular past, subject \textbf{ic}. Accusative \textbf{steorran} is both the object perceived and the actor of the supplied plain infinitive \textbf{feallan}. \textbf{Of heofone} gives origin; \textbf{oð eorðan} gives the endpoint, “as far as/\allowbreak{}until the earth”. The perception construction is already taught; supplying this infinitive does not assign a whole new productive verb system. A finite \textbf{fēoll} attached without a suitable clause construction misses the target.
% mdparagraph:P199:end

\end{minipage}\par\addvspace{4pt}

\par\begin{minipage}{\linewidth}\raggedright
% mdparagraph:P200:begin
\textbf{24.4.}
% mdparagraph:P200:end

% item: EX-VII-24-04
% key: KEY-EX-VII-24-04
% mdparagraph:P201:begin
\textbf{Prompt:} He that overcometh, the same shall be clothed (gescrȳdan) in white (hwīt) raiment.
% mdparagraph:P201:end

% mdparagraph:P202:begin
\textbf{Model:} Sē ðe oferwinð, se ilca bið mid hwītum rēafe gescrȳd.
% mdparagraph:P202:end

% mdparagraph:P203:begin
\textbf{Explanation:} The relative subject takes \textbf{oferwinð}, while \textbf{se ilca} resumes “the same one”. \textbf{Bið … gescrȳd} is the prospective passive; \textbf{mid hwītum rēafe} uses the evidenced clothing frame, with neuter dative agreement. English “in” does not require an Old English \textbf{on} phrase. The documented participle \textbf{gescrȳdd} is a valid alternative to selected \textbf{gescrȳd}. Preserve already taught \textbf{bið} without inventing a new quantity. [5]
% mdparagraph:P203:end

\end{minipage}\par\addvspace{4pt}

\par\begin{minipage}{\linewidth}\raggedright
% mdparagraph:P204:begin
\textbf{24.5.}
% mdparagraph:P204:end

% item: EX-VII-24-05
% key: KEY-EX-VII-24-05
% mdparagraph:P205:begin
\textbf{Prompt:} These things saith he that is holy.
% mdparagraph:P205:end

% mdparagraph:P206:begin
\textbf{Model:} Ðās ðing cwið sē ðe hālig is.
% mdparagraph:P206:end

% mdparagraph:P207:begin
\textbf{Explanation:} \textbf{Ðās ðing} is the fronted neuter plural object of singular \textbf{cwið}. The postposed subject is \textbf{sē ðe hālig is}, “he who is holy”. Inside that relative, \textbf{hālig} is a nominative predicate adjective, not an attributive weak form after \textbf{sē}. The sentence does not say that the things themselves speak.
% mdparagraph:P207:end

\end{minipage}\par\addvspace{4pt}

\par\begin{minipage}{\linewidth}\raggedright
% mdparagraph:P208:begin
\textbf{24.6.}
% mdparagraph:P208:end

% item: EX-VII-24-06
% key: KEY-EX-VII-24-06
% mdparagraph:P209:begin
\textbf{Prompt:} Fear him (ðone) that is able to destroy (fordōn) both soul and body in (on, 166) hell.
% mdparagraph:P209:end

% mdparagraph:P210:begin
\textbf{Model:} Ondrǣdað ðone ðe mæg fordōn ge sāwle ge līchaman on hęlle.
% mdparagraph:P210:end

% mdparagraph:P211:begin
\textbf{Explanation:} The model gives a plural command \textbf{ondrǣdað}; the prompt permits singular \textbf{ondrǣd} for one addressee. Accusative \textbf{ðone} is the person feared, represented as subject within \textbf{ðe mæg fordōn}. The modal takes plain \textbf{fordōn}; \textbf{ge sāwle ge līchaman} pairs its two accusative objects. \textbf{On hęlle} is the supplied locative phrase. Do not insert an independent “may” verb merely because the relative has modal ability.
% mdparagraph:P211:end

\end{minipage}\par\addvspace{4pt}

\par\begin{minipage}{\linewidth}\raggedright
% mdparagraph:P212:begin
\textbf{24.7.}
% mdparagraph:P212:end

% item: EX-VII-24-07
% key: KEY-EX-VII-24-07
% mdparagraph:P213:begin
\textbf{Prompt:} I am not worthy (wierðe) that thou enter under (172) my roof.
% mdparagraph:P213:end

% mdparagraph:P214:begin
\textbf{Model:} Ic ne eom wierðe ðæt ðū under mīnne hrōf ingā.
% mdparagraph:P214:end

% mdparagraph:P215:begin
\textbf{Explanation:} Uncontracted \textbf{ne eom} is fully supported. \textbf{Wierðe} introduces the bounded content frame with present subjunctive \textbf{ingā}, subject \textbf{ðū}. The entry phrase \textbf{under mīnne hrōf} has masculine accusative agreement. This is not a general rule that every adjective takes any subjunctive content clause. A negative-copula contraction is unnecessary, and indicative \textbf{ingǣst} would miss the assigned mood. [5]
% mdparagraph:P215:end

\end{minipage}\par\addvspace{4pt}

\par\begin{minipage}{\linewidth}\raggedright
% mdparagraph:P216:begin
\textbf{24.8.}
% mdparagraph:P216:end

% item: EX-VII-24-08
% key: KEY-EX-VII-24-08
% mdparagraph:P217:begin
\textbf{Prompt:} Great is your reward in heaven (plur.).
% mdparagraph:P217:end

% mdparagraph:P218:begin
\textbf{Model:} Miclu is ēowru mēd on heofonum.
% mdparagraph:P218:end

% mdparagraph:P219:begin
\textbf{Explanation:} Feminine \textbf{mēd} is the nominative subject. Fronted predicate \textbf{miclu} and possessive \textbf{ēowru} carry the selected short feminine nominative forms; the model uses plural “your”. Singular \textbf{ðīn} is defensible if the context addresses one person. \textbf{On heofonum} retains the requested plural heaven. The fronted adjective does not change the subject’s gender.
% mdparagraph:P219:end

\end{minipage}\par\addvspace{4pt}

\par\begin{minipage}{\linewidth}\raggedright
% mdparagraph:P220:begin
\textbf{24.9.}
% mdparagraph:P220:end

% item: EX-VII-24-09
% key: KEY-EX-VII-24-09
% mdparagraph:P221:begin
\textbf{Prompt:} If God so clothes (om. the) grass of the field, which (87. b) to-day is, and is (bið) to-morrow sent (āsęndan) into (on, 172) the oven, much (178) more he (om. shall) clothe you, O ye of little faith.
% mdparagraph:P221:end

% mdparagraph:P222:begin
\textbf{Model:} Gif God swā gescrȳt æceres gærs ðæt ðe tō-dæg is and tō-morgen bið on ofen āsęnd, micle mā hē gescrȳt ēow, gē lȳtles gelēafan.
% mdparagraph:P222:end

% mdparagraph:P223:begin
\textbf{Explanation:} The \textbf{gif} clause contains the relative \textbf{ðæt ðe tō-dæg is and tō-morgen bið … āsęnd}, referring to neuter \textbf{gærs}. Cook’s \textbf{87.b} requests the inflected demonstrative plus \textbf{ðe}: \textbf{ðæt} is neuter nominative singular as subject within that relative, although its antecedent \textbf{gærs} is accusative in the conditional clause. Bare \textbf{ðe} is grammatical but misses this stated strategy target. The finite \textbf{bið} and bare participle make a prospective passive; accusative \textbf{on ofen} marks its destination. The main clause is \textbf{micle mā hē gescrȳt ēow}, preserving Cook’s omission of an added “shall”. Instrumental adverb \textbf{micle}, \textbf{tō-morgen}, and address \textbf{gē lȳtles gelēafan} remain locally supplied. Trace the relative before returning to the main subject; a single unanalysed string of clauses loses that relation.
% mdparagraph:P223:end

\end{minipage}\par\addvspace{4pt}

\par\begin{minipage}{\linewidth}\raggedright
% mdparagraph:P224:begin
\textbf{24.10.}
% mdparagraph:P224:end

% item: EX-VII-24-10
% key: KEY-EX-VII-24-10
% mdparagraph:P225:begin
\textbf{Prompt:} Better is the soul than food (męte), and the body than its raiment.
% mdparagraph:P225:end

% mdparagraph:P226:begin
\textbf{Model:} Bętere is sēo sāwol ðonne męte and se līchama ðonne his rēaf.
% mdparagraph:P226:end

% mdparagraph:P227:begin
\textbf{Explanation:} The fronted comparative predicate \textbf{bętere} describes \textbf{sēo sāwol}. \textbf{Ðonne męte} gives the first comparison, and the second coordinate compares \textbf{se līchama} with \textbf{his rēaf}, with the corresponding “is better” understood. A full second predicate is \textbf{is bętera}, with masculine nominative \textbf{bętera} agreeing with \textbf{līchama}; feminine \textbf{bętere} belongs with \textbf{sāwol}. Preserve Cook’s hooked \textbf{męte} and the singular possessive relation in “its raiment”; neither comparison turns its second term into the subject of the first.
% mdparagraph:P227:end

\end{minipage}\par\addvspace{4pt}

\par\begin{minipage}{\linewidth}\raggedright
% mdparagraph:P228:begin
\textbf{24.11.}
% mdparagraph:P228:end

% item: EX-VII-24-11
% key: KEY-EX-VII-24-11
% mdparagraph:P229:begin
\textbf{Prompt:} Behold (behealdan) the birds of heaven; your heavenly Father feedeth (fēdan) them.
% mdparagraph:P229:end

% mdparagraph:P230:begin
\textbf{Model:} Behealdað heofones fuglas; ēower heofonlica fæder fēt hīe.
% mdparagraph:P230:end

% mdparagraph:P231:begin
\textbf{Explanation:} The model’s plural command \textbf{behealdað} corresponds to plural possessive \textbf{ēower}; a coherent singular command \textbf{beheald} and \textbf{ðīn} is defensible if one person is addressed. \textbf{Heofones fuglas} is the accusative object of beholding. Singular \textbf{ēower heofonlica fæder} is subject of \textbf{fēt}, and plural \textbf{hīe} refers back to the birds as object. The model chooses weak \textbf{heofonlica} after the possessive; strong masculine nominative \textbf{heofonlic} is also permitted. Distinguish third-singular \textbf{fēt} from a plural finite form.
% mdparagraph:P231:end

\end{minipage}\par\addvspace{4pt}

\par\begin{minipage}{\linewidth}\raggedright
% mdparagraph:P232:begin
\textbf{24.12.}
% mdparagraph:P232:end

% item: EX-VII-24-12
% key: KEY-EX-VII-24-12
% mdparagraph:P233:begin
\textbf{Prompt:} What shall (sculon) we eat, or what drink?
% mdparagraph:P233:end

% mdparagraph:P234:begin
\textbf{Model:} Hwæt sculon wē etan oððe hwæt drincan?
% mdparagraph:P234:end

% mdparagraph:P235:begin
\textbf{Explanation:} \textbf{Sculon} agrees with first-plural \textbf{wē} and takes plain \textbf{etan}; \textbf{drincan} shares the understood \textbf{sculon wē} in the coordinated question. Each \textbf{hwæt} is the interrogative object of its infinitive. Repeating the auxiliary and subject in the second question is grammatical; leaving them understood also preserves the sense. Do not make \textbf{drincan} a finite plural merely because the first clause has a plural subject.
% mdparagraph:P235:end

\end{minipage}\par\addvspace{4pt}

\reserve{7\baselineskip}
% mdparagraph:P236:begin
\section*{Cook Group VII. Exercise 25}
% mdparagraph:P236:end

\reserve{5\baselineskip}
\par\begin{minipage}{\linewidth}\raggedright
% mdparagraph:P237:begin
\textbf{25.1.}
% mdparagraph:P237:end

% item: EX-VII-25-01
% key: KEY-EX-VII-25-01
% mdparagraph:P238:begin
\textbf{Prompt:} He that (sē ðe) loveth not his brother, whom (87. b) he seeth, how can (mugan) he love God, whom he seeth not?
% mdparagraph:P238:end

% mdparagraph:P239:begin
\textbf{Model:} Sē ðe ne lufað his brōðor, ðone ðe hē gesiehð, hū mæg hē God lufian, ðone ðe hē ne gesiehð?
% mdparagraph:P239:end

% mdparagraph:P240:begin
\textbf{Explanation:} Initial \textbf{sē ðe} identifies the person who does not love his brother; the following \textbf{ðone ðe hē gesiehð} refers to that brother as object of seeing. The main question is \textbf{hū mæg hē God lufian}. Its final \textbf{ðone ðe hē ne gesiehð} refers to God, again as the object within the relative. Keep the two negations with their respective clauses and the plain infinitive after \textbf{mæg}. Neither the nearest noun alone nor translating every \textbf{ðe} identically establishes the full reference map.
% mdparagraph:P240:end

\end{minipage}\par\addvspace{4pt}

\par\begin{minipage}{\linewidth}\raggedright
% mdparagraph:P241:begin
\textbf{25.2.}
% mdparagraph:P241:end

% item: EX-VII-25-02
% key: KEY-EX-VII-25-02
% mdparagraph:P242:begin
\textbf{Prompt:} The time of my departure is nigh at hand.
% mdparagraph:P242:end

% mdparagraph:P243:begin
\textbf{Model:} Mīnre forðfōre tīd is nēah.
% mdparagraph:P243:end

% mdparagraph:P244:begin
\textbf{Explanation:} The subject is \textbf{tīd}, “time”; \textbf{mīnre forðfōre} is the feminine genitive phrase “of my departure”. Predicate \textbf{nēah} means “near, at hand”. The model does not require a separate verb corresponding to English “at hand”. A nominative possessive \textbf{mīn} inside that genitive phrase would miss its agreement.
% mdparagraph:P244:end

\end{minipage}\par\addvspace{4pt}

\par\begin{minipage}{\linewidth}\raggedright
% mdparagraph:P245:begin
\textbf{25.3.}
% mdparagraph:P245:end

% item: EX-VII-25-03
% key: KEY-EX-VII-25-03
% mdparagraph:P246:begin
\textbf{Prompt:} He did something (hwæthwega) of this sort, while his companions were going to town (burg).
% mdparagraph:P246:end

% mdparagraph:P247:begin
\textbf{Model:} Hē dyde hwæthwega swilces, ðā-hwīle-ðe his gefēran tō byrig ēodon.
% mdparagraph:P247:end

% mdparagraph:P248:begin
\textbf{Explanation:} The main clause uses singular \textbf{dyde} with \textbf{hē}; \textbf{hwæthwega swilces} is the supplied “something of this kind”. \textbf{Ðā-hwīle-ðe} introduces the simultaneous circumstance, with plural subject \textbf{his gefēran} and \textbf{ēodon}. The selected weak noun here is “companions”. \textbf{Tō byrig} uses the irregular dative of \textbf{burg}. Ordinary \textbf{ēodon} can render “were going”; a progressive auxiliary is not required.
% mdparagraph:P248:end

\end{minipage}\par\addvspace{4pt}

\par\begin{minipage}{\linewidth}\raggedright
% mdparagraph:P249:begin
\textbf{25.4.}
% mdparagraph:P249:end

% item: EX-VII-25-04
% key: KEY-EX-VII-25-04
% mdparagraph:P250:begin
\textbf{Prompt:} Be with the brethren.
% mdparagraph:P250:end

% mdparagraph:P251:begin
\textbf{Model:} Wes mid ðǣm brōðrum.
% mdparagraph:P251:end

% mdparagraph:P252:begin
\textbf{Explanation:} \textbf{Wes} is the singular command “be”; \textbf{mid ðǣm brōðrum} names the company in the dative plural. The prompt leaves addressee number unspecified, so a correctly taught plural command is valid if several people are addressed. The model does not turn “be with” into a new prefixed verb.
% mdparagraph:P252:end

\end{minipage}\par\addvspace{4pt}

\par\begin{minipage}{\linewidth}\raggedright
% mdparagraph:P253:begin
\textbf{25.5.}
% mdparagraph:P253:end

% item: EX-VII-25-05
% key: KEY-EX-VII-25-05
% mdparagraph:P254:begin
\textbf{Prompt:} Open the windows of the house.
% mdparagraph:P254:end

% mdparagraph:P255:begin
\textbf{Model:} Geopena ðæs hūses ēagðȳrlu.
% mdparagraph:P255:end

% mdparagraph:P256:begin
\textbf{Explanation:} The model selects singular \textbf{geopena}. \textbf{Ēagðȳrlu} is the selected neuter accusative plural, with genitive \textbf{ðæs hūses} identifying the house. A coherent plural command \textbf{geopeniað} is also valid if several addressees are meant. Preserve the noun’s syncopated plural rather than adding an ending to the complete singular.
% mdparagraph:P256:end

\end{minipage}\par\addvspace{4pt}

\par\begin{minipage}{\linewidth}\raggedright
% mdparagraph:P257:begin
\textbf{25.6.}
% mdparagraph:P257:end

% item: EX-VII-25-06
% key: KEY-EX-VII-25-06
% mdparagraph:P258:begin
\textbf{Prompt:} Love (plur.) your enemies, do (om. to) them well that hate you, that ye (om. may) be (196. f) children of your Father, which (87. b) is in heaven (plur.).
% mdparagraph:P258:end

% mdparagraph:P259:begin
\textbf{Model:} Lufiað ēowre fīend, dōð ðǣm wel ðe ēow hatiað, ðæt gē bēon ēowres Fæder bearn, sē ðe on heofonum is.
% mdparagraph:P259:end

% mdparagraph:P260:begin
\textbf{Explanation:} Cook expressly requires plural commands \textbf{lufiað, dōð}. \textbf{Ēowre fīend} is accusative object; \textbf{ðǣm … ðe ēow hatiað} is the dative group to whom good is done. Purpose \textbf{ðæt gē bēon … bearn} takes the taught subjunctive without an inserted modal. \textbf{Ēowres Fæder} is genitive with predicate \textbf{bearn}; \textbf{sē ðe on heofonum is} refers to the Father, not the children. Keep plural heaven and distinguish the dative recipients from the direct object of “hate”.
% mdparagraph:P260:end

\end{minipage}\par\addvspace{4pt}

\par\begin{minipage}{\linewidth}\raggedright
% mdparagraph:P261:begin
\textbf{25.7.}
% mdparagraph:P261:end

% item: EX-VII-25-07
% key: KEY-EX-VII-25-07
% mdparagraph:P262:begin
\textbf{Prompt:} Ye are my friends, if ye do (ind.) the things which (87. d) I command (bebēodan) you (verb last).
% mdparagraph:P262:end

% mdparagraph:P263:begin
\textbf{Model:} Gē sindon mīne frīend gif gē ðā ðing dōð ðe ic ēow bebēode.
% mdparagraph:P263:end

% mdparagraph:P264:begin
\textbf{Explanation:} \textbf{Sindon} and conditional indicative \textbf{dōð} agree with \textbf{gē}. \textbf{Mīne frīend} is the predicate nominative; \textbf{ðā ðing} is the object of doing. The relative \textbf{ðe ic ēow bebēode} has subject \textbf{ic}, dative recipient \textbf{ēow}, and the things as its understood object. We read Cook’s final-verb instruction as applying to the nearest named verb, \textbf{bebēode}, which is last in that relative. A defensible alternative places the entire relative beside \textbf{ðing} and leaves \textbf{dōð} last in its conditional. Do not reject either grammatical arrangement merely for not matching the model; retain the indicative and the stated relative-final target.
% mdparagraph:P264:end

\end{minipage}\par\addvspace{4pt}

\par\begin{minipage}{\linewidth}\raggedright
% mdparagraph:P265:begin
\textbf{25.8.}
% mdparagraph:P265:end

% item: EX-VII-25-08
% key: KEY-EX-VII-25-08
% mdparagraph:P266:begin
\textbf{Prompt:} The birds sang the sweetest song, winsome (acc.) to hear.
% mdparagraph:P266:end

% mdparagraph:P267:begin
\textbf{Model:} Ðā fuglas sungon ðone swētestan sang, wynsumne tō hīeranne.
% mdparagraph:P267:end

% mdparagraph:P268:begin
\textbf{Explanation:} Plural \textbf{ðā fuglas} takes strong past \textbf{sungon}. Its accusative object is \textbf{ðone swētestan sang}; \textbf{wynsumne} is the explicitly requested accusative masculine adjective agreeing with that song. \textbf{Tō hīeranne} is the taught inflected-infinitive frame, “to hear”. Nominative \textbf{wynsum} would miss Cook’s agreement control, even if an independently recast sentence could use it.
% mdparagraph:P268:end

\end{minipage}\par\addvspace{4pt}

\par\begin{minipage}{\linewidth}\raggedright
% mdparagraph:P269:begin
\textbf{25.9.}
% mdparagraph:P269:end

% item: EX-VII-25-09
% key: KEY-EX-VII-25-09
% mdparagraph:P270:begin
\textbf{Prompt:} The roof of the church is very high.
% mdparagraph:P270:end

% mdparagraph:P271:begin
\textbf{Model:} Ðǣre cirican hrōf is swīðe hēah.
% mdparagraph:P271:end

% mdparagraph:P272:begin
\textbf{Explanation:} \textbf{Ðǣre cirican} is a fronted genitive modifier within the subject phrase \textbf{ðǣre cirican hrōf}; the roof remains masculine nominative subject of \textbf{is}. Predicate \textbf{hēah} agrees with it and \textbf{swīðe} means “very”. Do not treat the initial church phrase as the subject or use a weak adjective merely because a genitive precedes the noun.
% mdparagraph:P272:end

\end{minipage}\par\addvspace{4pt}

\par\begin{minipage}{\linewidth}\raggedright
% mdparagraph:P273:begin
\textbf{25.10.}
% mdparagraph:P273:end

% item: EX-VII-25-10
% key: KEY-EX-VII-25-10
% mdparagraph:P274:begin
\textbf{Prompt:} There (ðā) was suddenly with the angel (om. a) multitude [number] of the heavenly host, praising (gen. plur.; cf. 189. 2) God, and thus saying.
% mdparagraph:P274:end

% mdparagraph:P275:begin
\textbf{Model:} Ðā wæs sęmninga mid ðǣm ęngle getæl heofonlices weorodes, hęriendra God and ðus cweðendra.
% mdparagraph:P275:end

% mdparagraph:P276:begin
\textbf{Explanation:} Cook supplies temporal/\allowbreak{}presentational \textbf{ðā}, not a new obligatory expletive corresponding to English “there”. Singular \textbf{getæl} is subject of \textbf{wæs}; \textbf{mid ðǣm ęngle} has the required dative. The local collective formula supplies \textbf{heofonlices weorodes} and substantive genitive plurals \textbf{hęriendra … cweðendra}, “of people praising … and saying”. Preserve their genitive plural control instead of changing them to nominatives to agree mechanically with singular \textbf{getæl}. Their complete participial paradigm is not a recall target.
% mdparagraph:P276:end

\end{minipage}\par\addvspace{4pt}

\par\begin{minipage}{\linewidth}\raggedright
% mdparagraph:P277:begin
\textbf{25.11.}
% mdparagraph:P277:end

% item: EX-VII-25-11
% key: KEY-EX-VII-25-11
% mdparagraph:P278:begin
\textbf{Prompt:} To God be (193) glory (wuldor) in (om. the) heights, and on earth peace (om. to) men of good will.
% mdparagraph:P278:end

% mdparagraph:P279:begin
\textbf{Model:} Gode sīe wuldor on hēanessum and on eorðan sibb mannum gōdes willan.
% mdparagraph:P279:end

% mdparagraph:P280:begin
\textbf{Explanation:} This is the supplied wish formula: dative recipient \textbf{Gode}, optative \textbf{sīe}, nominative \textbf{wuldor}. The coordinate wish has nominative \textbf{sibb} with \textbf{sīe} understood; \textbf{mannum} is dative without a separate \textbf{tō}, and \textbf{gōdes willan} is the bounded quality expression “of good will”. \textbf{On hēanessum} and \textbf{on eorðan} supply the two locations. Adding articles that Cook explicitly omits misses its control; treating \textbf{Gode} as a nominative subject changes the construction.
% mdparagraph:P280:end

\end{minipage}\par\addvspace{4pt}

\par\begin{minipage}{\linewidth}\raggedright
% mdparagraph:P281:begin
\textbf{25.12.}
% mdparagraph:P281:end

% item: EX-VII-25-12
% key: KEY-EX-VII-25-12
% mdparagraph:P282:begin
\textbf{Prompt:} He returned (ęfthweorfan) (om. by) the same way that he had come (see ǣr, adv.).
% mdparagraph:P282:end

% mdparagraph:P283:begin
\textbf{Model:} Hē ęfthwearf ðȳ ilcan wege ðe hē ǣr cōm.
% mdparagraph:P283:end

% mdparagraph:P284:begin
\textbf{Explanation:} Singular \textbf{ęfthwearf} means “returned”. \textbf{Ðȳ ilcan wege} is the taught route expression, “by the same way”, without an additional preposition. The relative \textbf{ðe hē ǣr cōm} establishes the earlier coming through \textbf{ǣr} and context. Plain preterite \textbf{cōm} therefore renders English “had come” without an obligatory compound pluperfect. The pronoun remains the person who returned; the reading later uses a similar route expression for the song.
% mdparagraph:P284:end

\end{minipage}\par\addvspace{4pt}

\clearpage
% mdparagraph:P285:begin
\section*{R. The bishop and the monk}
% mdparagraph:P285:end

\reserve{5\baselineskip}
\par\begin{minipage}{\linewidth}\raggedright
% mdparagraph:P286:begin
\textbf{R01. Complete translation.}
% mdparagraph:P286:end

% item: NEW-U13-R01
% key: KEY-NEW-U13-R01
% mdparagraph:P287:begin
The bishop sat in the house, and the monk was there. When the monk heard the song, he went to the bishop. The monk said, “I heard the song, which was very pleasant. Did you hear it?” The bishop said, “Go now, and bring the brothers here.” Then the monk went to the church and came back with the brothers. The bishop saw them and said, “You are my friends. I want to speak with you.” The monk sat with the brothers, and they heard the bishop’s words.
% mdparagraph:P287:end

% mdparagraph:P288:begin
\textbf{Target and alternatives:} Preserve the complete sequence, both speeches, the question and the command. “Lead the brothers here” is a valid literal rendering of \textbf{lǣd}; “words/\allowbreak{}speech” can render the final collective sense of \textbf{word}. The account does not give an explicit answer to the monk’s question, so “Yes, I heard it” must not be added. It is new teaching material, not an extra historical episode in Chad’s life.
% mdparagraph:P288:end

\end{minipage}\par\addvspace{4pt}

\par\begin{minipage}{\linewidth}\raggedright
% mdparagraph:P289:begin
\textbf{R02. Reference across the account.}
% mdparagraph:P289:end

% item: NEW-U13-R02
% key: KEY-NEW-U13-R02
% mdparagraph:P290:begin
In \textbf{ðā ēode hē}, \textbf{hē} is the monk. \textbf{Hine} in the monk’s question refers to masculine \textbf{sang}, so English uses “it”. In \textbf{geseah hīe}, the natural reference is the returned group, the monk and brothers. The final quotation’s \textbf{Ic} is the bishop, its speaker; \textbf{ēow} naturally addresses that group. A narrower reference of \textbf{hīe} or \textbf{ēow} to the brothers is also defensible: the wording does not force the exact membership of either plural group. Morphological gender and number constrain the reading, while context identifies the participants. Neither \textbf{hine} nor masculine \textbf{hē} must refer to a man merely by form.
% mdparagraph:P290:end

\end{minipage}\par\addvspace{4pt}

\par\begin{minipage}{\linewidth}\raggedright
% mdparagraph:P291:begin
\textbf{R03. Clauses and continued subjects.}
% mdparagraph:P291:end

% item: NEW-U13-R03
% key: KEY-NEW-U13-R03
% mdparagraph:P292:begin
The temporal clause is \textbf{[Ðā se munuc ðone sang hīerde]}, finite \textbf{hīerde}, subject \textbf{se munuc}. Its main clause is \textbf{[ðā ēode hē tō ðǣm bisceope]}, finite \textbf{ēode}, subject \textbf{hē}, the monk. The relative is \textbf{[ðe swīðe wynsum wæs]}, finite \textbf{wæs}; \textbf{ðe} represents its subject, the antecedent \textbf{sang}.
% mdparagraph:P292:end

% mdparagraph:P293:begin
One suitable coordination is \textbf{Ðā ēode se munuc tō ðǣre cirican, and cōm ęft mid ðǣm brōðrum}: \textbf{ēode} and \textbf{cōm} share \textbf{se munuc}. Another is \textbf{Se bisceop geseah hīe and cwæð}: \textbf{geseah} and \textbf{cwæð} share \textbf{se bisceop}. Give at least one actual pair and recover its subject. Do not carry that subject across a later explicitly different noun merely because another \textbf{and} occurs.
% mdparagraph:P293:end

\end{minipage}\par\addvspace{4pt}

\par\begin{minipage}{\linewidth}\raggedright
% mdparagraph:P294:begin
\textbf{R04. Sequence and final company.}
% mdparagraph:P294:end

% item: NEW-U13-R04
% key: KEY-NEW-U13-R04
% mdparagraph:P295:begin
The sequence is: the monk hears the song; he asks whether the bishop heard it; the bishop commands him to fetch the brothers; the monk returns with them; the bishop addresses the assembled group. The bishop, monk and brothers are present at the end. \textbf{Cōm ęft mid ðǣm brōðrum} brings the group back; \textbf{Se bisceop geseah hīe and cwæð} places the bishop with them; \textbf{Se munuc sæt mid ðǣm brōðrum} explicitly seats the monk among the brothers. The final \textbf{hīe hīerdon ðæs bisceopes word} presents the group listening. The text does not supply a number for these brothers or identify them as Chad’s seven.
% mdparagraph:P295:end

\end{minipage}\par\addvspace{4pt}

\reserve{8\baselineskip}
% mdparagraph:P296:begin
\section*{S. The Passing of Chad}
% mdparagraph:P296:end

\reserve{5\baselineskip}
\reserve{8\baselineskip}
% mdparagraph:P297:begin
\textbf{S01. Complete supported translation.}
% mdparagraph:P297:end

% item: NEW-U13-S01
% key: KEY-NEW-U13-S01
% mdparagraph:P298:begin
When, one day, he was doing something of this sort outside, and his companions were going to the town, to church, as they often did, and the bishop was alone in the oratory, occupied either in reading or in prayer, he suddenly heard—as he said afterwards—the sweetest and most beautiful sound of people singing and rejoicing descending from heaven to earth. He said that he first heard that sound and that song from the south-eastern part of the sky, that is, from the height of the winter sunrise; and from there they approached him little by little, until it reached the roof of the oratory in which the bishop was, and, entering, filled the whole place and surrounded it on all sides. He then directed his attention earnestly to what he heard. Then, after about half an hour, he heard the same song of joy rise from the roof of that same oratory and return upwards to heaven with indescribable sweetness, by the same way by which it had come before.
% mdparagraph:P298:end

% mdparagraph:P299:begin
He remained there for some time, wondering and marvelling; with an attentive mind he thought and considered what those things might mean. Then the bishop opened the window of the oratory and made a sign with his hand, as was his custom if anyone was outside, that he should come in to him. He immediately went in to him. He—the bishop—said to him: “Go quickly to the church and tell our seven brothers to come here to me; and you too be present with us.”
% mdparagraph:P299:end

% mdparagraph:P300:begin
\textbf{Target and alternatives:} Preserve all actions, the sound’s full route, the interval, the signal and the final summons. “His companions went/\allowbreak{}had gone” is a defensible contextual rendering of \textbf{ēodon}; the Old English preterite does not alone distinguish those English aspects. “Voice” can render \textbf{stefne}. The plural sound-and-song description gives way to singular masculine \textbf{sang}, naturally English “it”. “The whole place” makes the supplied singular \textbf{eall} explicit; “filled everything there” is also defensible. The interpretive \textbf{hwæt … bēon sceolden} may be rendered “what these things could signify/\allowbreak{}were to mean”, without an obligatory assertion of future obligation. \textbf{Mid wes} permits “you too be there/\allowbreak{}be present”; “with us” makes its implied company explicit. The scene’s first listener is Owini, not the bishop. The communal church differs from the bishop’s local oratory. Do not add the brothers’ arrival, the exhortation or Chad’s death: they lie beyond the assigned endpoint. [6]
% mdparagraph:P300:end

\reserve{5\baselineskip}
\par\begin{minipage}{\linewidth}\raggedright
% mdparagraph:P301:begin
\textbf{S02. The long opening sentence.}
% mdparagraph:P301:end

% item: NEW-U13-S02
% key: KEY-NEW-U13-S02
% mdparagraph:P302:begin
The major structure is \textbf{[When he one day was doing something of this kind outside, and his companions were going to town, to church [as they often did], and the bishop alone in the oratory was occupied either in reading or prayer], [then he suddenly heard [as he afterwards said] the sweetest and most beautiful sound … descend from heaven to earth].} The subordinate comparisons and parenthetical reporting interrupt a simple linear sequence.
% mdparagraph:P302:end

% mdparagraph:P303:begin
The subjects are: \textbf{dyde — hē}, Owini; \textbf{ēodon — his gefēran}, the companions; \textbf{dydon — hīe}, the companions; \textbf{wæs — se bisceop}, Chad; main \textbf{gehīerde — hē}, Owini; parenthetical \textbf{sægde — hē}, also Owini. The first \textbf{Ðā} introduces the temporal setting, the second \textbf{ðā} is an internal temporal adverb “then”, and the \textbf{ðā} immediately before \textbf{gehīerde} resumes the main hearing event. The further \textbf{ðā} before \textbf{swētestan stefne} is an accusative feminine article, not another clause marker. \textbf{Āstīgan} is an infinitive whose understood actor is the accusative sound; it is not another finite past verb. Do not make the nearby bishop the subject of hearing merely because his name occurs last.
% mdparagraph:P303:end

\end{minipage}\par\addvspace{4pt}

\par\begin{minipage}{\linewidth}\raggedright
% mdparagraph:P304:begin
\textbf{S03. The sound’s path.}
% mdparagraph:P304:end

% item: NEW-U13-S03
% key: KEY-NEW-U13-S03
% mdparagraph:P305:begin
Both \textbf{hē becōm} and \textbf{ðe hē ǣr cōm} refer to the masculine song; Cook’s note 152.9a explicitly says \textbf{Se sang}. It descends from heaven to earth, is first heard from the south-eastern sky at the height of the winter sunrise, approaches little by little, reaches the oratory’s roof, enters, fills and surrounds the place, and later rises from that same roof and returns upwards to heaven by the same route. The opening plural \textbf{nēalǣcton} treats the sound and song together; the later singular pronoun picks up \textbf{sang}. English “it” is appropriate. Saying that Owini flies to the roof or heaven confuses grammatical gender with human reference.
% mdparagraph:P305:end

\end{minipage}\par\addvspace{4pt}

\par\begin{minipage}{\linewidth}\raggedright
% mdparagraph:P306:begin
\textbf{S04. Owini and Chad.}
% mdparagraph:P306:end

% item: NEW-U13-S04
% key: KEY-NEW-U13-S04
% mdparagraph:P307:begin
Owini hears and reports the sound, directs his attention to it, hears its ascent and return, remains wondering and marvelling, thinks about its meaning, and enters when signalled. Chad is alone in the oratory, opens its window, makes the customary hand signal, and tells Owini to fetch the brothers. The exact \textbf{Cwæð hē, se bisceop, him tō} makes the quoted speaker explicit: appositive \textbf{se bisceop} identifies \textbf{hē} as Chad. Here \textbf{him} is Owini, the addressee; the earlier \textbf{hē} of hearing is Owini, while the intervening motion pronouns can refer to the song. Grouping every masculine pronoun under one actor would destroy these distinctions.
% mdparagraph:P307:end

\end{minipage}\par\addvspace{4pt}

\par\begin{minipage}{\linewidth}\raggedright
% mdparagraph:P308:begin
\textbf{S05. The stranded preposition.}
% mdparagraph:P308:end

% item: NEW-U13-S05
% key: KEY-NEW-U13-S05
% mdparagraph:P309:begin
The whole phrase is “to the roof of the oratory in which the bishop was”: \textbf{tō ðæcean ðǣre ciricean ðe se bisceop in wæs}. The antecedent of \textbf{ðe} is the oratory/\allowbreak{}church, \textbf{ciricean}, not the roof. Cook note 152.11 says that \textbf{ðe} is governed by \textbf{in}, although that preposition stands later in the relative. \textbf{Se bisceop} is its nominative subject and \textbf{wæs} its finite verb. The particle’s unchanged spelling does not exhibit a dative ending; the prepositional relationship supplies the function. The roof is the endpoint of the song’s arrival, while the bishop is inside the oratory.
% mdparagraph:P309:end

\end{minipage}\par\addvspace{4pt}

\par\begin{minipage}{\linewidth}\raggedright
% mdparagraph:P310:begin
\textbf{S06. Fronting from the content clause.}
% mdparagraph:P310:end

% item: NEW-U13-S06
% key: KEY-NEW-U13-S06
% mdparagraph:P311:begin
The two fronted accusatives are \textbf{ðā stefne} and \textbf{ðone sang}. Their governing verb is embedded \textbf{gehīerde}, as Cook note 152.8 states. \textbf{Hē cwæð} is the reporting clause; \textbf{ðæt hē ǣrest gehīerde …} is its content. A useful map is \textbf{[Fronted objects: ðā stefne and ðone sang] [Report: hē cwæð [Content: ðæt hē ǣrest gehīerde \_\_\_ fram ēastsūðdǣle heofones …]].} The gap notation is explanatory, not an addition to the source.
% mdparagraph:P311:end

% mdparagraph:P312:begin
The meaning is “He said that he first heard that sound and song …”, not “He spoke that sound and song”. In the dependent clause, finite \textbf{gehīerde} precedes the source phrase \textbf{fram ēastsūðdǣle heofones}; the actual source therefore also supplies a concrete example of a dependent verb that is not last. Do not erase this order to fit a universal rule.
% mdparagraph:P312:end

\end{minipage}\par\addvspace{4pt}

\par\begin{minipage}{\linewidth}\raggedright
% mdparagraph:P313:begin
\textbf{S07. The summons.}
% mdparagraph:P313:end

% item: NEW-U13-S07
% key: KEY-NEW-U13-S07
% mdparagraph:P314:begin
Chad is the speaker and Owini the singular addressee. \textbf{Ūre} is the inclusive communal “our”, referring to their brotherhood; \textbf{mē} points to Chad, the speaking bishop; \textbf{ðū} singles out Owini. The command is to have \textbf{ūre seofon brōðor}, “our seven brothers”, come \textbf{hider tō mē}, here to the bishop. \textbf{And ðū ēac swilce mid wes} explicitly includes Owini as well; it does not send him away from the gathering.
% mdparagraph:P314:end

% mdparagraph:P315:begin
The fetch command \textbf{Gang hraðe tō ciricean} directs Owini to the communal church. Cook note 152.5 distinguishes that church, where the companions have gone, from the oratory in which Chad is waiting. The exact source’s church spellings are supported forms, not a request to reconstruct an unassigned proper-name or dialect paradigm.
% mdparagraph:P315:end

\end{minipage}\par\addvspace{4pt}

\par\begin{minipage}{\linewidth}\raggedright
% mdparagraph:P316:begin
\textbf{S08. Two actual clause orders.}
% mdparagraph:P316:end

% item: NEW-U13-S08
% key: KEY-NEW-U13-S08
% mdparagraph:P317:begin
One suitable main clause is \textbf{Ðā ontȳnde se bisceop ðæt ēagðȳrel ðǣre ciricean}. Its first constituent is adverb \textbf{ðā}; finite \textbf{ontȳnde} follows, then subject \textbf{se bisceop} and the complete window object phrase. This illustrates the common fronted-constituent plus finite-verb pattern.
% mdparagraph:P317:end

% mdparagraph:P318:begin
One suitable dependent clause is \textbf{gif hwilc mann ūte wǣre}. Its introductory conjunction is \textbf{gif}, its first following constituent the subject phrase \textbf{hwilc mann}, its finite verb \textbf{wǣre}; that verb is late, after adverb \textbf{ūte}. Other correctly identified clauses are acceptable. For example, \textbf{ðæt hē ǣrest gehīerde fram ēastsūðdǣle heofones} has introductory \textbf{ðæt}, subject \textbf{hē}, adverb \textbf{ǣrest}, finite \textbf{gehīerde}, then the source phrase; it demonstrates a permitted nonfinal dependent verb. Distinguish a clause introducer from the constituents following it. Preserve the source wording and justify the actual relation to a tendency; do not manufacture a compulsory pattern.
% mdparagraph:P318:end

\end{minipage}\par\addvspace{4pt}

\reserve{8\baselineskip}
% mdparagraph:P319:begin
\section*{Source and editorial notes}
% mdparagraph:P319:end

\begingroup\fontsize{10}{12.8}\selectfont
% mdparagraph:P320:begin
\textbf{[1]} Cook, \emph{Exercises in Old English} (1895), Group VII, printed pp. 11–13, one-based PDF pp. 19–21. Original prompts are preserved; all answer models are new teaching material. “Omit” directions and bracketed controls belong to the original prompts. No model string is the sole test of valid Old English.
% mdparagraph:P320:end

% mdparagraph:P321:begin
\textbf{[2]} The nominal systems follow the selected Cook lexemes and the complete declension rules documented in the learner chapter. Kinship and syncopating forms require their stated exceptions. Comparative \textbf{swētra} follows the weak system and the consonant simplification checked in Wright §259.3. Forms derived from those rules are distinguished from individual lexical attestations.
% mdparagraph:P321:end

% mdparagraph:P322:begin
\textbf{[3]} Cook’s selected strong and weak verb rules, lexical entries and the chapter’s separately assigned command, infinitive and participle systems. For \textbf{forceorfan}, consult the strong III framework and checked lexical evidence; original reference \textbf{134} points to \textbf{munan} in this edition. Its cause is unproved. Preserve the prompt and the documented correction rather than silently changing its numeral.
% mdparagraph:P322:end

% mdparagraph:P323:begin
\textbf{[4]} Fulk, \emph{An Introductory Grammar of Old English}, §78, qualifies main- and dependent-clause ordering descriptions. Cook’s exact \textbf{Æfter ðon hē underðēodde} is on printed p. 153, editorial line 16 after the printed line-15 marker. It is a separately supplied fragment, outside the assigned continuous reading.
% mdparagraph:P323:end

% mdparagraph:P324:begin
\textbf{[5]} The dictionary evidence for \textbf{beforan} positioning, \textbf{under} roof entry, static \textbf{wið} plus accusative, and \textbf{gescrȳdan} with \textbf{mid} plus dative is distinct from the morphology of those nouns and verbs. Cook §172.1 and the \textbf{niman} entry, printed p. 307, jointly support the bounded declared composition in 23.9; they do not attest that whole sentence as a quotation. \textbf{Tō-morgen} is supplied from Wright, printed p. 294, one-based PDF p. 314.
% mdparagraph:P324:end

% mdparagraph:P325:begin
\textbf{[6]} Cook, \emph{A First Book in Old English}, second edition (1895), printed 152.3–153.9 partial, one-based PDF pp. 168–169. Only multiples of five are printed marginal line numbers; intervening line positions are editorial counts. The reading ends at \textbf{and ðū ēac swilce mid wes.} Notes 152.5, 152.8, 152.9a, 152.11 and 152.12 establish respectively the church/\allowbreak{}oratory distinction, the embedded accusatives, the song reference, the stranded preposition and singular \textbf{eall}. All 24 original notes remain in the learner chapter. The preceding context is supplied in English; the following exhortation and death are unassigned.
% mdparagraph:P325:end

% mdparagraph:P326:begin
Original apparatus questions are retained as source notes, not added graded tasks. The key answers every assigned learner task, including the two complete readings; supplied etymological questions do not create additional recall requirements.
% mdparagraph:P326:end

\endgroup
\end{document}
